\documentclass[11pt,a4paper]{article}

\usepackage[utf8]{inputenc}
\usepackage[T1]{fontenc}
\usepackage{geometry}
\geometry{margin=1in}

% Paragraph spacing
\setlength{\parskip}{0.6\baselineskip plus 2pt minus 1pt}
\setlength{\parindent}{0pt}

\usepackage{mathtools,amsmath,amssymb,amsthm}
\usepackage{booktabs}
\usepackage[hidelinks]{hyperref}
\usepackage{enumitem}
\usepackage{url}
\usepackage{cleveref}

% Theorem environments
\theoremstyle{plain}
\newtheorem{theorem}{Theorem}[section]
\newtheorem{lemma}[theorem]{Lemma}
\newtheorem{corollary}[theorem]{Corollary}
\newtheorem{proposition}[theorem]{Proposition}
\theoremstyle{definition}
\newtheorem{definition}[theorem]{Definition}
\newtheorem{axiom}[theorem]{Axiom}
\newtheorem{assumption}[theorem]{Assumption}
\newtheorem{example}[theorem]{Example}
\theoremstyle{remark}
\newtheorem{remark}[theorem]{Remark}

\newcommand{\defeq}{\vcentcolon=}
\newcommand{\nf}{\mathrm{nf}}
\newcommand{\apply}{\mathrm{apply}}
\newcommand{\Fed}{\mathrm{Fed}}
% Allow line breaks after underscores in Coq names set in \texttt.
\renewcommand{\_}{\textunderscore\allowbreak}
\setlength{\emergencystretch}{3em}

\title{Normalization Confluence in Federated Registry Networks}

\author{
Dayna Blackwell\\
\small\href{mailto:dayna@blackwell-systems.com}{dayna@blackwell-systems.com}\\
\small Licensed under CC-BY-4.0.}

\date{4 October 2026 (Version 2)}

\begin{document}
\maketitle

\begin{center}
\small Version 2. See Appendix (Errata and corrections) for changes from version 1.
\end{center}

\begin{abstract}
\setlength{\parskip}{0.6\baselineskip plus 2pt minus 1pt}
Normalization confluence~\cite{blackwell2026nc} guarantees coordination-free
convergence within a single registry: under well-founded compensation (WFC)
and compensation commutativity (CC), all processors consuming the same events
reach the same valid state regardless of application order.  We extend this
guarantee to \emph{federated} environments where multiple registries, each
with independent invariants and compensation, are connected by
\emph{registry morphisms} encoding cross-organizational semantic constraints.

Our key construction treats a federation of registries connected by morphisms
as a registry, with morphism consistency conditions as additional invariants
and morphism repair as additional compensation.  Repair confluence then
composes freely: on any acyclic network whose morphisms preserve validity (M1),
with a source-determined (R1), validity-preserving (R2) \emph{resolution
operator} at each multi-source target, one round of federated normalization
reaches the unique federated normal form, from any state.  The core mechanism
is an \emph{authority argument}: the sources' states determine every target's
shared component.  Event confluence does not compose for free.  A target event
can read a shared value that morphism repair overwrites, so component
convergence and M1 do not imply federated convergence (we give counterexamples
satisfying every component condition).  We identify two local, per-edge checks,
cross-registry compensation commutativity (C1) and repaired compensation
commutativity (C2), and prove that they suffice for convergence of every event
interleaving when events are applied to normalized states, and that the same two
equations, evaluated at the witnesses a run can reach, are exactly necessary and
sufficient.

We prove the structural conditions necessary by counterexample: without
acyclicity, contradictory cyclic morphisms cause compensation to cycle;
without validity preservation, repair and local compensation oscillate; without
source determinacy, the shared normal form depends on the interleaving.  On
ordered shared domains with monotone repair, cycles are admitted: the repair
normal form is the least fixed point (federally valid when bottom is valid), and
event interleavings converge exactly when independent events commute after
re-normalization, which per-target C1 and C2 over the repair images imply.  The
federated results are mechanized in Coq/Rocq, axiom-free.  For
self-containment, we restate the single-registry model and convergence results
of~\cite{blackwell2026nc} alongside a verification calculus for CC.
\end{abstract}

\medskip\noindent\textbf{Keywords:} Stream Processing, Normalization
Confluence, Semantic Consistency, Rewrite Systems, Newman's Lemma,
Compensation Commutativity, Federated Systems, Registry Morphisms,
Distributed Systems

\clearpage

%========================================
\section{Introduction}
\label{sec:intro}

Coordination-free convergence---the ability of distributed processors to agree
on a result without explicit synchronization---is among the central goals of
distributed systems design.  The literature has identified two structural
regimes under which it is achievable:

\begin{enumerate}[leftmargin=*,nosep]
\item \textbf{Operation commutativity} (CRDTs~\cite{shapiro2011crdt}):
  Operations are designed to commute algebraically.  Any application order
  produces the same state.
\item \textbf{Invariant confluence}
  (I-confluence~\cite{bailis2014coordination}): Operations individually
  preserve application invariants under all orderings.  No repair is needed
  because validity is never violated.
\end{enumerate}

Both regimes achieve convergence by \emph{preventing} inconsistency.  CRDTs
prevent it through algebraic design; I-confluence prevents it through
invariant-preserving operation selection.  Neither accommodates operations that
are non-commutative \emph{and} may violate invariants---yet such operations
are pervasive in practice.  An order processing pipeline, for example, must
handle events (orders, payments, inventory adjustments, shipments) that do not
commute and that individually violate business rules (no shipment without
payment, no negative inventory).

We formalize the missing piece as a \emph{registry}: any authoritative
definition of validity external to the stream processor.  Registries already
exist in practice under various names---policy engines, constraint services,
schema validators, compliance controllers---but their role in convergence has
not been characterized formally.  Our model abstracts these as a tuple of
invariant predicates and a compensation operator.

This paper presents a unified treatment of a third regime, \textbf{normalization
confluence}~\cite{blackwell2026nc}, and extends it to multi-registry
environments.  In normalization confluence, operations may be non-commutative
and may violate invariants.  But a
compensation operator normalizes states to validity after each operation, and
the key structural condition is that these \emph{compensated results}
commute---even though the operations themselves do not.  Under this condition,
together with well-founded compensation, Newman's Lemma guarantees that all
causality-respecting processing orders reach the same valid state.

The three regimes form a hierarchy of progressively weaker requirements:
\begin{enumerate}[leftmargin=*,nosep]
\item Operation commutativity $\implies$ compensation commutativity
  $\implies$ convergence.
\item Invariant confluence $\implies$ no compensation needed $\implies$
  convergence.
\item Compensation commutativity alone $\implies$ convergence~\cite{blackwell2026nc}.
\end{enumerate}
Normalization confluence occupies the design space between CRDT commutativity
(which requires operations to commute) and I-confluence (which requires
operations to preserve invariants).  It permits operations that satisfy
neither, provided their compensated results agree.

Modern distributed systems, however, span organizational boundaries: a supply
chain links manufacturers and suppliers, each with its own registry.  We show
that normalization confluence extends naturally to these \emph{federated}
settings.  Multiple registries connected by \emph{morphisms}---structured
mappings encoding cross-organizational constraints---form a single federated
registry.  Termination and the repair normal form need nothing beyond the
component conditions and validity preservation of the morphisms (on acyclic
networks, with \emph{resolution operators}, source-determined and
validity-preserving merges, at multi-source targets).  Convergence of event
interleavings needs more: a target event may read a shared value that a
morphism overwrites, and then component convergence does not carry over.  Two
local checks per edge, C1 and C2, close the gap, and their reachable forms are
exactly necessary and sufficient.

\paragraph{Relationship to prior work.}
This paper extends~\cite{blackwell2026nc}, which establishes the
single-registry governance rewrite system, the core convergence result
(WFC~+~CC $\Rightarrow$ unique normal forms), necessity and complexity results,
and a verification calculus for CC based on invariant footprints and
decomposable repair.  The present paper adds the federated extension
(\Cref{sec:federation}): registry morphisms with partial authority, the
federated construction, the authority argument, the federated convergence
theorem under the checks C1 and C2 with its exact converse and necessity
counterexamples, and its generalization to multi-source acyclic networks via
resolution operators (\Cref{sec:multisource}).  For
self-containment, we restate the
single-registry model and include all proofs needed to support the federated
result; readers familiar with~\cite{blackwell2026nc} may proceed directly to
\Cref{sec:federation}.

\subsection{Convergence via Normalization}

Convergence is conventionally treated as a property of \emph{operations}:
operations must be designed to commute, or designed to preserve invariants.
Our result reframes convergence as a property of the \emph{normalization
rewrite system}: the system that interleaves event application with invariant
repair.  The operations themselves need no special algebraic structure.  What
matters is that the rewrite system---encompassing both application and
compensation---is terminating and confluent.  This shifts the design burden
from operation algebra to compensation design.

\subsection{Contributions}

The single-registry results (items 1--6) are from~\cite{blackwell2026nc},
restated here for self-containment.  The federated extension (items 7--10) is new.

\begin{enumerate}[leftmargin=*,nosep]
\item \textbf{Normalization confluence regime} (\Cref{sec:model}): A
  governance rewrite system on configurations $(\sigma, B)$ pairing a state
  with a pending-event buffer, with rules for both event application and
  compensation.  Nondeterminism arises from the choice of which enabled event
  to apply next.
\item \textbf{Convergence theorem} (\Cref{sec:convergence}): Proof that
  processors consuming the same events converge to the same valid state,
  regardless of application order, under well-founded compensation (WFC) and
  compensation commutativity (CC).
\item \textbf{Necessity} (\Cref{sec:necessity}):
  Counterexamples demonstrating that without UBC, compensation depth grows
  without uniform bound (necessary for complexity guarantees), and without CC,
  different orders reach different valid states (necessary for convergence).
\item \textbf{Complexity analysis} (\Cref{sec:complexity}): Proof that
  registry governance adds $O(\log n)$ overhead per event under UBC,
  preserving practical scalability.
\item \textbf{Three-regime characterization} (\Cref{sec:discussion}):
  Together with CRDTs and invariant confluence, a partition of the
  coordination-free design space into operation commutativity, invariant
  confluence, and normalization confluence.
\item \textbf{Verification calculus} (\Cref{sec:calculus}): Sufficient
  conditions---invariant footprints and decomposable repair---that reduce the
  global CC obligation to local, per-event-pair checks.  Product composition
  enables modular verification of independent registries.
\item \textbf{Federated convergence} (\Cref{sec:federation}): Extension to
  tree-shaped multi-registry networks connected by morphisms.  Registry
  morphisms with partial authority model cross-organizational constraints.
  Federated termination and a unique repair normal form are derived from
  network structure via an authority argument, with only morphism validity
  preservation (M1) needed.  Event convergence is not derived: component CC
  and M1 do not suffice (counterexamples).  Under two per-edge checks, C1 and C2,
  every event interleaving applied to normalized states converges
  (\Cref{thm:fed-convergence}); the same equations at reachable witnesses are
  the exact condition.  Necessity of acyclicity, M1, C1 and C2 is shown by
  counterexample.
\item \textbf{Multi-source resolution} (\Cref{sec:multisource}): Removal of the
  tree restriction.  A target with several incoming morphisms carries a
  \emph{resolution operator} generalizing the authority function; under source
  determinacy and validity preservation (R1--R2), the repair normal form is
  unique on any acyclic network, and event convergence holds under C1 and C2
  (\Cref{thm:resolved-convergence}), with the tree theorem as the
  single-source special case.  R1, R2, C1 and C2 are finite conditions,
  checkable by enumeration, so a resolved network is certified as WFC and CC
  are.
\item \textbf{Monotone cycles} (\Cref{sec:monotone}): When shared domains are
  ordered lattices and repair is monotone, acyclicity is unnecessary for the
  repair normal form: it is the least fixed point (Knaster--Tarski), reached by
  every fair chaotic iteration under the ascending chain condition, and
  federally valid when every component is valid at bottom
  (\Cref{thm:monotone-cycles}).  Event interleavings on such a network converge
  exactly when independent events commute after re-normalization (the global
  condition GC), which per-target C1 and C2 over the repair images imply.
  Acyclicity and monotonicity are two independent routes to a unique repair
  normal form.
\item \textbf{Compositionality} (\Cref{sec:compositionality}): A convex,
  internally convergent sub-federation collapses to a single \emph{effective
  registry} (\Cref{thm:collapse}), with internal convergence now including C1
  and C2 (or GC); the outer network and the collapsed one have the same repair
  normal form and the same runs.  This gives modular and hierarchical
  verification, generalizing product composition from independent registries to
  sub-diagrams.  This result is proved on paper; it is not yet mechanized.
\item \textbf{Machine-checked results} (\Cref{sec:mechanization},
  \Cref{rem:fed-mechanized}): The single-registry Convergence Theorem and the
  federated results above, including every counterexample and the exact
  converses, are mechanized in Coq/Rocq, axiom-free (\texttt{Print
  Assumptions} reports ``Closed under the global context'' for every theorem in
  the gate), with concrete instances witnessing that the hypotheses are jointly
  satisfiable.
\end{enumerate}

%========================================
\section{Related Work}
\label{sec:related}

\paragraph{Conflict-Free Replicated Data Types.}
CRDTs~\cite{shapiro2011crdt} guarantee convergence through algebraic
structure: replicas that apply the same set of commutative, associative,
idempotent operations converge to the same state.  Our setting is strictly more
general in one dimension: we do not require operations to commute, only that
their \emph{compensated results} commute.  This permits operations that
individually violate invariants, which CRDTs cannot accommodate.  Conversely,
CRDTs require no finiteness or boundedness assumptions.  Neither approach
dominates: CRDTs are preferable when operations can be designed to commute,
while registry governance is preferable when business rules constrain valid
states in ways that preclude commutativity.

\paragraph{CALM and Coordination Avoidance.}
The CALM theorem~\cite{hellerstein2010declarative,ameloot2013relational}
establishes that monotone programs can achieve consistency without coordination.
Bailis et al.~\cite{bailis2014coordination} extend this to identify
\emph{invariant confluence}: a set of operations is invariant-confluent if all
possible orderings preserve application invariants.  When invariant confluence
holds, no compensation is needed.  Our framework addresses the complement:
operations that are \emph{not} invariant-confluent.  We prove that
compensation restores convergence under bounded depth and commutativity.
Together, invariant confluence and our result partition the design space
(\Cref{sec:design-space}).

\paragraph{Invariant-Repair Systems.}
Indigo~\cite{balegas2015indigo} and IPA~\cite{balegas2019ipa} also restore
application invariants over weakly consistent replicas, and Indigo in particular
targets the cross-object constraints that arise across service boundaries.  They
do so by a different route: static analysis of invariant-violating operation
pairs, to which they attach application-specific conflict-resolution rules or
reservation (escrow) protocols so that every reachable state stays
invariant-preserving, in the invariant-confluence tradition.  We do not
preselect operations to be invariant-preserving.  We admit operations that
genuinely violate invariants and prove, via Newman's Lemma under compensation
commutativity, that repair reconciles all orderings to the same valid normal
form; the federated extension carries this across registry morphisms.  The
distinction is one of method: reservation and invariant confluence versus
confluence of a compensating rewrite system.

\paragraph{Abstract Rewrite Systems.}
Newman's Lemma~\cite{newman1942} states that a terminating abstract rewrite
system is confluent if and only if it is locally confluent.  The lemma has been
applied extensively in term rewriting~\cite{baader1998term} and lambda
calculus.  Our contribution is its application to a rewrite system with
\emph{genuine nondeterminism}: the choice of which event to apply next from a
buffer of enabled events.  Local confluence is nontrivial because different
application orders produce different intermediate states; our compensation
commutativity axiom is the condition that closes the resulting critical pairs.

\paragraph{Rewriting Logic and Coherence.}
Meseguer's rewriting logic~\cite{meseguer1992} models concurrent systems as
rewrite theories, and coherence checking (the Church-Rosser and coherence
properties of rewrite theories modulo equations~\cite{duran2012coherence},
supported by tools such as Maude) is the closest methodological ancestor of our
build-time convergence check.  We instantiate this lineage for one specific
rewrite system: events interleaved with invariant compensation, where
compensation commutativity plays the role of the coherence condition that closes
critical pairs.

\paragraph{Self-Stabilization.}
Self-stabilizing systems~\cite{dijkstra1974} converge to a legitimate state from
an arbitrary starting configuration.  Normalization confluence shares the
converge-to-valid premise but strengthens the conclusion to agreement: not
merely that each replica reaches some valid state, but that all replicas, across
a federation, reach the \emph{same} valid state under reordering.

\paragraph{Eventual Consistency.}
Classical eventual consistency~\cite{vogels2009eventually} guarantees that
replicas converge given sufficient time without new updates.  It does not
address whether converged states satisfy application-level invariants---only
that replicas agree on the same value.  Our work guarantees convergence to
\emph{valid} states under continuous event arrival.

\paragraph{Stream Processing Theory.}
Formal models of stream processing include synchronous
dataflow~\cite{lee1987synchronous}, Kahn process
networks~\cite{kahn1974semantics}, and the windowed computation models
underlying modern engines~\cite{akidau2015dataflow}.  These focus on
operational semantics---what a stream processor computes---rather than semantic
validity of computed states.  Our registry layer sits above any of these
operational models.

\paragraph{Saga Pattern and Compensating Transactions.}
The saga pattern~\cite{garcia1987sagas} decomposes long-lived transactions into
sequences of local transactions with compensating actions for rollback.  Our
compensation operator generalizes this: rather than undoing individual
operations, it repairs arbitrary invariant violations in the current state.
The convergence guarantee we prove---that different interleaving orders of
application and compensation reach the same valid state---has no analogue in
the saga literature.

\paragraph{Mechanized Convergence Proofs.}
Gomes et al.~\cite{gomes2017verifying} give a machine-checked proof of strong
eventual consistency for CRDTs in Isabelle/HOL, network model included.  Their
result covers the commutative case; the mechanization this work builds on
(\Cref{sec:mechanization}) covers convergence under non-commutative,
invariant-violating operations reconciled by compensation.  Both rest on an
assumption-free, machine-checked convergence theorem rather than a paper proof.

%========================================
\section{Registry-Governed Stream Model}
\label{sec:model}

\subsection{Registries and Semantic Projection}

\begin{definition}[Registry]
\label{def:registry}
A \emph{registry} is a tuple $R = (\Sigma_R, I_R, V_R, \rho_R)$ where:
\begin{itemize}[leftmargin=*,nosep]
\item $\Sigma_R$ is a set of \emph{semantic states},
\item $I_R = \{\psi_1, \ldots, \psi_k\}$ is a finite set of \emph{invariant
  predicates} $\psi_i: \Sigma_R \to \{\top, \bot\}$,
\item $V_R: \Sigma_R \to \{\top, \bot\}$ is the \emph{validation function}
  defined by $V_R(\sigma) = \bigwedge_{i=1}^k \psi_i(\sigma)$,
\item $\rho_R: \Sigma_R \to \Sigma_R$ is a deterministic \emph{compensation
  operator} such that if $V_R(\sigma) = \top$, then
  $\rho_R(\sigma) = \sigma$.
\end{itemize}
\end{definition}

\begin{definition}[Semantic Projection]
\label{def:projection}
A registry is defined over a semantic state space $\Sigma_R$ together with
a \emph{projection} $\pi: \Sigma_{\mathrm{raw}} \to \Sigma_R$ from
unbounded raw states to semantic states.  All invariants, validation, and
compensation are defined on $\Sigma_R$.  Two raw states $s_1, s_2$ are
\emph{registry-indistinguishable} if $\pi(s_1) = \pi(s_2)$.
The application function $\mathrm{apply}: \mathcal{E} \times \Sigma_R \to \Sigma_R$
is understood as ``apply then project'': given a raw-level application
$\mathrm{apply}_{\mathrm{raw}}: \mathcal{E} \times \Sigma_{\mathrm{raw}} \to \Sigma_{\mathrm{raw}}$,
the semantic application is
$\mathrm{apply}(e, \sigma) = \pi(\mathrm{apply}_{\mathrm{raw}}(e, s))$ for
any $s$ with $\pi(s) = \sigma$.  Well-definedness requires that the projected
result is independent of the representative~$s$.
\end{definition}

\begin{remark}[Scope of Finiteness]
\label{rem:scope}
Finiteness is required of the \emph{semantic} state space $\Sigma_R$, not the
raw data.  An order-processing registry tracking status across states
\emph{pending}, \emph{paid}, \emph{shipped}, \emph{delivered},
\emph{cancelled} has $|\Sigma_R| = 5$ regardless of how many bytes the
underlying order record occupies.  More generally, any domain whose logic is
captured by $k$ boolean predicates has $|\Sigma_R| \leq 2^k$: the semantic
states are the equivalence classes under the invariant vector
$(\psi_1(\sigma), \ldots, \psi_k(\sigma))$.  Business domains are
overwhelmingly of this form.  The finiteness assumption applies to these
semantic equivalence classes, not to raw data representations.
\end{remark}

\begin{definition}[Registry Equivalence]
\label{def:equiv}
Two states $\sigma_1, \sigma_2 \in \Sigma_R$ are \emph{registry equivalent},
written $\sigma_1 \approx_R \sigma_2$, if they satisfy the same invariants:
$\forall \psi \in I_R: \psi(\sigma_1) = \psi(\sigma_2)$.
\end{definition}

\begin{lemma}
\label{lem:equiv-relation}
Registry equivalence $\approx_R$ is an equivalence relation on~$\Sigma_R$.
\end{lemma}
\begin{proof}
Each condition $\psi(\sigma_1) = \psi(\sigma_2)$ is reflexive, symmetric, and
transitive.  A conjunction of equivalence relations is an equivalence relation.
\end{proof}

\subsection{Well-Founded and Uniformly Bounded Compensation}

Termination and complexity require two properties of the compensation
operator, which we state separately because they play different roles.

\begin{axiom}[Well-Founded Compensation (WFC)]
\label{ax:wfc}
There exists a \emph{compensation measure}
$\Phi: \Sigma_R \to \mathbb{N}$, defined over semantic states (after
projection), such that every compensation step strictly
decreases $\Phi$: if $V_R(\sigma) = \bot$, then
$\Phi(\rho_R(\sigma)) < \Phi(\sigma)$.
\end{axiom}

WFC guarantees that compensation terminates from any state: iterated
application of~$\rho_R$ reaches a valid state in finitely many steps.

\begin{axiom}[Uniformly Bounded Compensation (UBC)]
\label{ax:ubc}
The compensation measure $\Phi$ from \Cref{ax:wfc} is uniformly bounded:
there exists $M \in \mathbb{N}$ such that $\Phi(\sigma) \leq M$ for all
$\sigma \in \Sigma_R$.
\end{axiom}

UBC strengthens WFC by bounding compensation depth \emph{uniformly} across
all states.  UBC does not require finiteness of $\Sigma_R$; bounded
compensation height alone suffices.  WFC alone suffices for termination and
for defining iterated
compensation~$\rho_R^*$.  UBC is additionally needed for the uniform
step bound in the termination lemma and the constant per-event compensation
cost in the complexity analysis.

\begin{lemma}[Finite State Spaces Imply UBC]
\label{lem:finite-implies-ubc}
If $|\Sigma_R| < \infty$ and WFC holds, then UBC holds with
$M = \max_{\sigma \in \Sigma_R} \Phi(\sigma)$.
\end{lemma}
\begin{proof}
The maximum exists because $\Sigma_R$ is finite.
\end{proof}

\begin{remark}
WFC is the load-bearing condition for termination; UBC is the load-bearing
condition for uniform step bounds and per-event complexity.  Finiteness of
$\Sigma_R$ is a convenient sufficient condition for both, but UBC does not
require full finiteness of $\Sigma_R$; it requires only bounded compensation
height.  The counterexample
in \Cref{sec:necessity} satisfies WFC (compensation terminates for every fixed
event set) but violates UBC (compensation depth grows without bound as event
sets grow).
\end{remark}

\subsection{Events and Streams}

\begin{definition}[Event Stream]
\label{def:stream}
An \emph{event stream} $\mathcal{S} = \langle e_1, e_2, \ldots \rangle$ is a
(possibly infinite) sequence of events drawn from event space $\mathcal{E}$.
Each event $e_i$ carries a unique identifier $\mathrm{id}(e_i)$, a timestamp
$\tau(e_i)$, and a set of causal dependencies
$\mathrm{deps}(e_i) \subseteq \{e_1, \ldots, e_{i-1}\}$.  An
\emph{application function}
$\mathrm{apply}: \mathcal{E} \times \Sigma_R \to \Sigma_R$ maps each event
and current state to a new state.
\end{definition}

\begin{definition}[Causality-Respecting Order]
\label{def:causal-order}
An ordering $\pi$ of a finite event set $E$ \emph{respects causality} if
whenever $e_i \in \mathrm{deps}(e_j)$, event $e_i$ precedes $e_j$
under~$\pi$.  Multiple causality-respecting orders may exist for the same
event set.
\end{definition}

\subsection{Governance Rewrite System}

We now define the rewrite system that is the core technical device of this
paper.  Unlike the approach of fixing a canonical event order---which would make
agreement trivial by determinism---our rewrite system admits genuine
nondeterminism in the choice of which event to apply next.

\begin{assumption}[Causal Closure]
\label{as:closure}
The received event set $E$ is causally closed: for every $e \in E$,
$\mathrm{deps}(e) \subseteq E$.
\end{assumption}

Under causal closure, every dependency of every event in~$E$ is either still
pending in~$B$ or has already been applied (removed from~$B$).  This makes the
enabled predicate below well-defined.

\begin{definition}[Configuration]
\label{def:config}
A \emph{configuration} is a pair $(\sigma, B)$ where $\sigma \in \Sigma_R$ is
a semantic state and $B \subseteq E$ is a set of events that have been
received but not yet applied.  An event $e \in B$ is \emph{enabled} if
$\mathrm{deps}(e) \cap B = \varnothing$ (all causal dependencies of $e$ have
already been applied).
\end{definition}

\begin{definition}[Governance Rewrite System]
\label{def:grs}
Fix a finite set of received events~$E$ and initial state~$\sigma_0$.  The
\emph{governance rewrite system} $\mathcal{G} = (\mathcal{C}, \to)$ operates
on configurations $(\sigma, B)$ with $\sigma \in \Sigma_R$ and
$B \subseteq E$.  The rewrite relation $\to$ is the union of:
\begin{enumerate}[leftmargin=*,nosep]
\item \textbf{Apply step:}
  $(\sigma, B) \to (\mathrm{apply}(e, \sigma), B \setminus \{e\})$ whenever
  $e \in B$ is enabled.
\item \textbf{Compensation step:}
  $(\sigma, B) \to (\rho_R(\sigma), B)$ whenever $V_R(\sigma) = \bot$.
\end{enumerate}
A configuration is in \emph{normal form} if $B = \varnothing$ and
$V_R(\sigma) = \top$.
\end{definition}

The nondeterminism is genuine: when multiple events are enabled and the state
is invalid, a configuration may step via different apply rules (choosing
different events) or via the compensation rule.  Confluence of this system
means that the choice does not affect the final result.

\subsection{Compensation Commutativity}

Local confluence of $\mathcal{G}$ is nontrivial.  Two critical pair types
arise: (1)~two apply steps choosing different events, and (2)~an apply step
versus a compensation step.  The following axiom ensures both types close.

\begin{definition}[Iterated Compensation]
\label{def:rhostar}
Under WFC (\Cref{ax:wfc}), define $\rho_R^*(\sigma)$ to be the unique state
obtained by iterating $\rho_R$ from~$\sigma$ until reaching validity:
$\rho_R^*(\sigma) = \rho_R^{m}(\sigma)$ for the least $m$ such that
$V_R(\rho_R^{m}(\sigma)) = \top$.  Existence and finiteness of~$m$ follow from
WFC; uniqueness follows from determinism of~$\rho_R$.
\end{definition}

\begin{axiom}[Compensation Commutativity (CC)]
\label{ax:cc}
The registry $R$ satisfies:
\begin{enumerate}[leftmargin=*,nosep]
\item[\textbf{(CC1)}] \textbf{Order independence.}\; For any
  $\sigma \in \Sigma_R$ and events $e_1, e_2$ that are causally independent
  and simultaneously enabled in some configuration $(\sigma, B)$:
  \[
  \rho_R^*(\mathrm{apply}(e_2,\, \rho_R^*(\mathrm{apply}(e_1, \sigma))))
  \;=\;
  \rho_R^*(\mathrm{apply}(e_1,\, \rho_R^*(\mathrm{apply}(e_2, \sigma)))).
  \]
\item[\textbf{(CC2)}] \textbf{Compensation absorption.}\; For any
  $\sigma \in \Sigma_R$ with $V_R(\sigma) = \bot$ and any enabled event~$e$:
  \[
  \rho_R^*(\mathrm{apply}(e, \sigma))
  \;=\;
  \rho_R^*(\mathrm{apply}(e, \rho_R(\sigma))).
  \]
\end{enumerate}
\end{axiom}

\begin{remark}[CRDT Commutativity]
\label{rem:crdt-comparison}
CRDT commutativity requires
$\mathrm{apply}(e_1, \mathrm{apply}(e_2, \sigma)) =
\mathrm{apply}(e_2, \mathrm{apply}(e_1, \sigma))$ for all~$\sigma$.  This
implies CC1 (take $\rho_R^* = \mathrm{id}$ when all states are valid) and
trivially satisfies CC2 (when all states are valid, compensation never fires).
Compensation commutativity is strictly weaker: it permits operations that
individually violate invariants, requiring only that the \emph{repaired}
results commute.  The gap between operation commutativity and compensation
commutativity is precisely the space our result occupies.
\end{remark}

\begin{remark}[Compensation Absorption]
\label{rem:absorption}
Condition CC2 states that compensating an invalid state before applying an
event does not change the final compensated result.  Intuitively, the
repair that compensation performs ``would have been performed anyway'' during
the post-application compensation phase.  CC2 holds whenever compensation
repairs invariant violations in a way that is determined by the violation
pattern rather than the path taken to reach the violation.  By induction on
the number of compensation steps, CC2 generalizes:
$\rho_R^*(\mathrm{apply}(e, \sigma)) =
\rho_R^*(\mathrm{apply}(e, \rho_R^*(\sigma)))$.
In rewrite-theoretic terms, CC2 ensures that the normalizer is stable under
interleaving with event application: normalizing ``early'' (before an event)
versus ``late'' (after) does not change the final normal form.
\end{remark}

\subsection{Stream Processors}

\begin{definition}[Stream Processor]
\label{def:processor}
A \emph{stream processor} $P$ for stream $\mathcal{S}$ under registry $R$ is a
process that, at each time $t$, has received some subset of events
$E_P(t) \subseteq \mathcal{S}$ and maintains state $\Psi_P(t) \in \Sigma_R$.
We require:
\begin{enumerate}[leftmargin=*,nosep]
\item \textbf{Eventual delivery}: For every event $e_i \in \mathcal{S}$, there
  exists a time $t_i$ such that $e_i \in E_P(t)$ for all $t \geq t_i$.
\item \textbf{Causality-respecting application}: The processor applies events
  in some order that respects causality (\Cref{def:causal-order}), buffering
  events whose causal dependencies have not yet been received.  The processor
  is free to choose any causality-respecting order.
\item \textbf{Incremental compensation}: After applying each event, the
  processor applies $\rho_R$ until validity is restored.
\end{enumerate}
The processor's computation on event set $E_P(t)$ is a reduction sequence in
the governance rewrite system $\mathcal{G}$ starting from
$(\sigma_0, E_P(t))$.
\end{definition}

Different processors may apply the same events in different
causality-respecting orders.  The convergence theorem guarantees they
nonetheless arrive at the same valid state.

%========================================
\section{Worked Example: Order Fulfillment}
\label{sec:example}

We instantiate the framework on an order fulfillment pipeline, verify the
structural conditions, and trace two processor orderings to the same valid
state.

\subsection{Registry Definition}

An order has a \emph{stage} in $\{\textit{pending}, \textit{approved},
\textit{held}\}$ and a \emph{balance} in $\{0, 1, 2, 3\}$.  The semantic
state space is $\Sigma_R = \{\textit{pending}, \textit{approved},
\textit{held}\} \times \{0, 1, 2, 3\}$, with $|\Sigma_R| = 12$ and initial
state $\sigma_0 = (\textit{pending}, 0)$.

A single invariant governs fulfillment:
\[
\psi(s, b) \;=\; \top \;\iff\; (s = \textit{approved} \implies b \geq 2).
\]
Approval requires a balance of at least~2.  The only invalid states are
$(\textit{approved}, 0)$ and $(\textit{approved}, 1)$.

\paragraph{Compensation.}
When an order is approved with insufficient balance, it is placed on hold:
\[
\rho(\textit{approved}, b) = (\textit{held}, b) \quad \text{for } b < 2.
\]
All other states are valid, and $\rho$ acts as the identity.  The
\textit{held} stage preserves the intent to approve: the system records that
approval was requested but cannot yet proceed.  This is a deliberate design
choice---a naive alternative that simply cancels the approval (reverting to
\textit{pending}) would violate CC, as we discuss below.

\paragraph{Measures.}
$\Phi(s, b) = 1$ if $s = \textit{approved}$ and $b < 2$, otherwise $0$.
WFC holds ($\rho$ strictly decreases $\Phi$ on invalid states), and UBC holds
with $M = 1$.

\subsection{Events}

Two causally independent events:
\begin{itemize}[leftmargin=*,nosep]
\item $e_{\mathrm{credit}}$: Adds 1 to the balance (capped at~3).  If the
  order is on hold and the new balance reaches~2, the order is automatically
  approved:
  \[
  \mathrm{apply}(e_{\mathrm{credit}}, (s, b)) =
  \begin{cases}
  (\textit{approved}, \min(b{+}1, 3)) & \text{if } s = \textit{held} \text{ and } \min(b{+}1,3) \geq 2, \\
  (s, \min(b{+}1, 3)) & \text{otherwise.}
  \end{cases}
  \]
\item $e_{\mathrm{approve}}$: Sets the stage to \textit{approved}:
  $\mathrm{apply}(e_{\mathrm{approve}}, (s, b)) = (\textit{approved}, b)$.
\end{itemize}

\subsection{Verifying CC}

\paragraph{CC1 (Order Independence).}
We verify CC1 for $(e_{\mathrm{approve}}, e_{\mathrm{credit}})$ at two
representative starting states.

\emph{From $(\textit{pending}, 0)$}:
\begin{align*}
&\text{Path 1 ($e_{\mathrm{approve}}$ first):} \\
&\quad (\textit{pending}, 0) \xrightarrow{e_{\mathrm{approve}}}
  (\textit{approved}, 0) \xrightarrow{\rho^*} (\textit{held}, 0)
  \xrightarrow{e_{\mathrm{credit}}} (\textit{held}, 1).
  \quad\text{Result: } (\textit{held}, 1). \\[4pt]
&\text{Path 2 ($e_{\mathrm{credit}}$ first):} \\
&\quad (\textit{pending}, 0) \xrightarrow{e_{\mathrm{credit}}}
  (\textit{pending}, 1) \xrightarrow{e_{\mathrm{approve}}}
  (\textit{approved}, 1) \xrightarrow{\rho^*} (\textit{held}, 1).
  \quad\text{Result: } (\textit{held}, 1). \;\checkmark
\end{align*}

In Path~1, approval fires first, violates the invariant (balance~0 $< 2$),
and compensation places the order on hold.  The subsequent credit increments
the balance but does not yet reach the threshold.  In Path~2, credit arrives
first (no violation), then approval violates the invariant and compensation
again produces the held state.  Both paths reach $(\textit{held}, 1)$.

\emph{From $(\textit{pending}, 1)$}:
\begin{align*}
&\text{Path 1:}\;
  (\textit{pending}, 1) \xrightarrow{e_{\mathrm{approve}}}
  (\textit{approved}, 1) \xrightarrow{\rho^*} (\textit{held}, 1)
  \xrightarrow{e_{\mathrm{credit}}} (\textit{approved}, 2).
  \quad\text{Result: } (\textit{approved}, 2). \\
&\text{Path 2:}\;
  (\textit{pending}, 1) \xrightarrow{e_{\mathrm{credit}}}
  (\textit{pending}, 2) \xrightarrow{e_{\mathrm{approve}}}
  (\textit{approved}, 2).
  \quad\text{Result: } (\textit{approved}, 2). \;\checkmark
\end{align*}

In Path~1, approval again violates the invariant; compensation holds; the
subsequent credit reaches balance~2, triggering auto-approval.  In Path~2, the
credit arrives first, and approval at balance~2 is valid.  Both reach
$(\textit{approved}, 2)$: the order is approved with sufficient funds.

Since $|\Sigma_R| = 12$, CC1 can be verified exhaustively for all valid
starting states.  All cases follow the same pattern: compensation preserves
the approval intent in the held state, and subsequent credits eventually
satisfy the balance threshold.

\paragraph{CC2 (Compensation Absorption).}
From the invalid state $(\textit{approved}, 0)$ with event
$e_{\mathrm{credit}}$:
\begin{align*}
\text{LHS:}\;\; & \rho^*(\mathrm{apply}(e_{\mathrm{credit}},\,
  (\textit{approved}, 0)))
  = \rho^*((\textit{approved}, 1))
  = (\textit{held}, 1). \\
\text{RHS:}\;\; & \rho^*(\mathrm{apply}(e_{\mathrm{credit}},\,
  \rho(\textit{approved}, 0)))
  = \rho^*(\mathrm{apply}(e_{\mathrm{credit}},\,
  (\textit{held}, 0)))
  = \rho^*((\textit{held}, 1))
  = (\textit{held}, 1). \;\checkmark
\end{align*}

Compensating before the credit (putting the order on hold first) produces the
same result as applying the credit to the invalid state and then compensating.
CC2 holds because the held state and the approved-but-invalid state respond
identically to subsequent credits after compensation.

\subsection{Two Processors, Three Events}

Consider the stream $\mathcal{S} = \langle e_{\mathrm{approve}},
e_{\mathrm{credit}}, e_{\mathrm{credit}} \rangle$ (approval followed by two
credits).  Two processors $P_1, P_2$ receive the events in different orders:

\medskip\noindent
\textbf{Processor $P_1$} (approve, credit, credit):
\[
(\textit{pending}, 0)
\xrightarrow{e_{\mathrm{approve}}} (\textit{approved}, 0)
\xrightarrow{\rho^*} (\textit{held}, 0)
\xrightarrow{e_{\mathrm{credit}}} (\textit{held}, 1)
\xrightarrow{e_{\mathrm{credit}}} (\textit{approved}, 2).
\]

\noindent
\textbf{Processor $P_2$} (credit, credit, approve):
\[
(\textit{pending}, 0)
\xrightarrow{e_{\mathrm{credit}}} (\textit{pending}, 1)
\xrightarrow{e_{\mathrm{credit}}} (\textit{pending}, 2)
\xrightarrow{e_{\mathrm{approve}}} (\textit{approved}, 2).
\]

\noindent Both processors arrive at $(\textit{approved}, 2)$.  $P_1$ encountered
an invariant violation (approval without sufficient balance) and passed
through the held state; $P_2$ never violated the invariant.  Compensation
commutativity guarantees agreement regardless of the path.

\subsection{Why Compensation Design Matters}

The held stage is essential.  A naive compensation that simply reverts
approval to \textit{pending}---i.e.,
$\rho(\textit{approved}, b) = (\textit{pending}, b)$---would violate CC1.
From $(\textit{pending}, 0)$, approval-then-credit would produce
$(\textit{pending}, 1)$ (intent to approve lost), while credit-then-approval
would produce $(\textit{held}, 1)$ or a different compensated state depending
on the design.  The held stage preserves the approval intent so that
subsequent credits can trigger auto-approval, ensuring that different orderings
converge.

This illustrates the practical content of the CC conditions: they constrain
\emph{compensation design}, not operation design.  The sufficient condition
used here is \emph{state-determined compensation}: $\rho^*$ depends only on
which invariants are violated (an approved order with low balance is always
placed on hold), making the post-compensation state independent of how the
violation was reached.

%========================================
\section{Convergence Theorem}
\label{sec:convergence}

\subsection{Termination}

\begin{lemma}[Termination]
\label{lem:termination}
Under WFC (\Cref{ax:wfc}), the governance rewrite system $\mathcal{G}$ is
terminating: every reduction sequence from $(\sigma_0, E)$ is finite.
Under UBC (\Cref{ax:ubc}), the total number of steps is bounded by
$|E| + (|E| + 1) \cdot M$.
\end{lemma}
\begin{proof}
Define the measure $\mu(\sigma, B) = (|B|, \Phi(\sigma))$ taking values in
$\mathbb{N} \times \mathbb{N}$ under the lexicographic order, which is
well-founded.

An apply step sends $(\sigma, B)$ to $(\mathrm{apply}(e, \sigma), B \setminus \{e\})$.
The first component decreases: $|B \setminus \{e\}| = |B| - 1$.  Therefore
$\mu$ strictly decreases regardless of the second component.

A compensation step sends $(\sigma, B)$ to $(\rho_R(\sigma), B)$.  The first
component is unchanged; the second strictly decreases by WFC.
Therefore $\mu$ strictly decreases.

Since $\mu$ maps into a well-ordered set and strictly decreases at every step,
every reduction sequence is finite.  Under UBC, $\Phi(\sigma) \leq M$ for all
$\sigma$, so at most $M$ compensation steps can occur between apply steps
(including before the first apply), giving the uniform bound
$|E| + (|E| + 1) \cdot M$.
\end{proof}

\subsection{Local Confluence}

\begin{lemma}[Local Confluence]
\label{lem:local-confluence}
Under CC (\Cref{ax:cc}), the governance rewrite system $\mathcal{G}$ is
locally confluent: if $(\sigma, B) \to X_1$ and $(\sigma, B) \to X_2$,
then there exists a configuration~$Z$ with $X_1 \to^* Z$ and
$X_2 \to^* Z$.
\end{lemma}
\begin{proof}
Three cases arise from the two rewrite rules.

\medskip\noindent\textbf{Case 1: Two apply steps.}\;
$(\sigma, B) \to (\mathrm{apply}(e_1, \sigma), B \setminus \{e_1\})
= X_1$ and
$(\sigma, B) \to (\mathrm{apply}(e_2, \sigma), B \setminus \{e_2\})
= X_2$, where $e_1 \neq e_2$ are both enabled.

Since $e_1$ and $e_2$ are both enabled, their dependencies are disjoint
from~$B$.  By \Cref{as:closure}, $e_1, e_2 \in E$, and since $e_1 \in B$ we
have $e_1 \notin \mathrm{deps}(e_2)$ (otherwise
$\mathrm{deps}(e_2) \cap B \neq \varnothing$, contradicting enabledness of
$e_2$), and symmetrically $e_2 \notin \mathrm{deps}(e_1)$.  Therefore $e_1$
and $e_2$ are causally independent.

From $X_1$: compensate to validity, obtaining
$(\rho_R^*(\mathrm{apply}(e_1, \sigma)),\, B \setminus \{e_1\})$; then
apply~$e_2$ (still enabled since its dependencies remain outside
$B \setminus \{e_1\}$) and compensate, reaching
\[
Z_1 = \bigl(\rho_R^*(\mathrm{apply}(e_2,\,
  \rho_R^*(\mathrm{apply}(e_1, \sigma)))),\;
  B \setminus \{e_1, e_2\}\bigr).
\]
From $X_2$: symmetrically, reach
\[
Z_2 = \bigl(\rho_R^*(\mathrm{apply}(e_1,\,
  \rho_R^*(\mathrm{apply}(e_2, \sigma)))),\;
  B \setminus \{e_1, e_2\}\bigr).
\]
By CC1, the state components of $Z_1$ and $Z_2$ are equal.  Take $Z = Z_1 = Z_2$.

\medskip\noindent\textbf{Case 2: Apply vs.\ compensation.}\;
$V_R(\sigma) = \bot$ and $e$ is enabled in~$B$.
$(\sigma, B) \to (\mathrm{apply}(e, \sigma), B \setminus \{e\}) = X_1$
and
$(\sigma, B) \to (\rho_R(\sigma), B) = X_2$.

From $X_1$: compensate to validity, reaching
$(\rho_R^*(\mathrm{apply}(e, \sigma)),\, B \setminus \{e\})$.

From $X_2 = (\rho_R(\sigma), B)$: apply~$e$ (still enabled, since
compensation does not change~$B$) to get
$(\mathrm{apply}(e, \rho_R(\sigma)),\, B \setminus \{e\})$; then
compensate to validity, reaching
$(\rho_R^*(\mathrm{apply}(e, \rho_R(\sigma))),\, B \setminus \{e\})$.

By CC2, $\rho_R^*(\mathrm{apply}(e, \sigma)) =
\rho_R^*(\mathrm{apply}(e, \rho_R(\sigma)))$.  The buffer components
are equal.  Take $Z = (\rho_R^*(\mathrm{apply}(e, \sigma)),\, B \setminus
\{e\})$.

\medskip\noindent\textbf{Case 3: Two compensation steps.}\;
Both rewrites apply $\rho_R$ to the same state, producing the same
successor $(\rho_R(\sigma), B)$.  Take $Z = (\rho_R(\sigma), B)$.
\end{proof}

\subsection{Main Result}

The following results are from~\cite{blackwell2026nc}; proofs are included
for self-containment.

\begin{corollary}[Unique Normal Forms]
\label{cor:unique-nf}
Under WFC and CC, every configuration $(\sigma_0, E)$ has a unique normal
form.
\end{corollary}
\begin{proof}
By \Cref{lem:termination}, $\mathcal{G}$ is terminating.  By
\Cref{lem:local-confluence}, $\mathcal{G}$ is locally confluent.  By Newman's
Lemma~\cite{newman1942}, $\mathcal{G}$ is globally confluent.  A terminating,
confluent rewrite system has unique normal forms.
\end{proof}

\begin{remark}[Machine-checked]
\label{sec:mechanization}
The single-registry base results above (\cite{blackwell2026nc}) are mechanized
in Coq/Rocq (checked with Rocq~9.3): Newman's Lemma is proved from scratch, and
the Governance Rewrite System's termination (via the lexicographic measure
$(|B|, \Phi(\sigma))$) and local confluence (the three critical-pair cases,
closed by CC1, CC2, and determinism of~$\rho_R$) yield confluence and unique
normal forms.  For both the abstract theorem and a concrete instantiation,
\texttt{Print Assumptions} reports ``Closed under the global context'': no
axioms, no admitted lemmas.  A concrete state-determined-compensation registry
discharges WFC and CC with genuine nondeterminism, so the hypotheses are jointly
satisfiable and the theorem is not vacuous.  The soundness of the reference
implementation's certification strategy is mechanized as well---footprint
disjointness implies order-independent commutation (the compositional CC path),
and a repair strictly decreasing a potential is strongly normalizing (WFC)---each
axiom-free.  The monotone-cycles result (\Cref{thm:monotone-cycles}) is mechanized
in its constructive finite-lattice core: Kleene iteration of a monotone operator
from~$\bot$ is an ascending chain bounded by every fixed point, and once it
stabilizes the stable value is the (unique) least fixed point~$\bar{s}^{\star}$,
again axiom-free with a concrete witness.  The chaotic-iteration convergence---that
every fair \emph{asynchronous} component-wise schedule from~$\bot$, not only the
synchronous iteration, reaches~$\bar{s}^{\star}$---is mechanized as well: reframing
a chaotic schedule as a productive-update rewrite relation, below~$\bar{s}^{\star}$
a monotone update strictly increases a rank, so the relation is strongly
normalizing and its unique normal form is~$\bar{s}^{\star}$; hence any productive
schedule reaches the same fixed point (order-independence), again axiom-free with a
concrete witness.  The single-registry development is furthermore paired with two
runnable checkers extracted from it, each proven axiom-free: one certifies that a
compiled machine's emitted step tables commute and stay in range, and one
recomputes convergence directly from the machine's rules by evaluating their
expression trees (trusting neither the implementation's enumeration nor its
normalization).  Extracted from the proof rather than reimplemented by hand, they
act as independent differential oracles for the reference implementation: a defect
in its hand-written verification cannot make a non-convergent machine pass either
checker.  Scripts are in the \texttt{coq/} directory of the artifact.
\end{remark}

\begin{theorem}[Stream Convergence]
\label{thm:convergence}
Let $R$ be a registry satisfying WFC (\Cref{ax:wfc}) and CC
(\Cref{ax:cc}), and let $P_1, P_2$ be stream processors
(\Cref{def:processor}) consuming the same event stream~$\mathcal{S}$.
Then:
\begin{enumerate}[leftmargin=*,nosep,label=(\alph*)]
\item \textbf{Validity}: For each processor, after it completes the
  incremental compensation phase for any applied event, its state satisfies
  $V_R(\Psi_{P_j}(t)) = \top$.
\item \textbf{Order independence}: The state of each processor depends only on
  its received event set, not on the order in which events were applied.
\item \textbf{Agreement}: Whenever the processors have received the same
  events, they are in the same state: if $E_{P_1}(t) = E_{P_2}(t)$, then
  $\Psi_{P_1}(t) = \Psi_{P_2}(t)$.
\end{enumerate}
Combined with eventual delivery, disagreements are transient: any event causing
processors to differ will eventually be received by both, restoring agreement.
\end{theorem}

\begin{proof}
We prove each part.

\medskip\noindent\textbf{Part (a): Validity.}\;
By \Cref{lem:termination}, compensation terminates.  By
\Cref{def:processor}, the processor compensates after each event
application until validity is restored.  Therefore after completing the
compensation phase for any applied event,
$V_R(\Psi_{P_j}(t)) = \top$.

\medskip\noindent\textbf{Part (b): Order independence.}\;
Each processor's computation on event set $E = E_{P_j}(t)$ is a reduction
sequence in the governance rewrite system from $(\sigma_0, E)$ to a normal
form $(\sigma^*, \varnothing)$.  Different causality-respecting orders
correspond to different reduction sequences.  By \Cref{cor:unique-nf}, all
reduction sequences reach the same normal form.  Therefore the final
state~$\sigma^*$ depends only on $E$, not on the order.

\medskip\noindent\textbf{Part (c): Agreement.}\;
If $E_{P_1}(t) = E_{P_2}(t) = E$, then both processors' computations are
reduction sequences from $(\sigma_0, E)$.  By \Cref{cor:unique-nf}, both
reach the same normal form $(\sigma^*, \varnothing)$.  Therefore
$\Psi_{P_1}(t) = \Psi_{P_2}(t) = \sigma^*$.
\end{proof}

\begin{remark}[Semantic Projection]
\label{rem:set-function}
The theorem implies that the mapping $F(E) = \nf(\sigma_0, E)$---the unique
normal form of the governance rewrite system starting from $(\sigma_0, E)$---is
a well-defined function of the event \emph{set}~$E$, independent of causal
linearization.  This is the Church--Rosser property for our operational
semantics.
\end{remark}

\begin{corollary}[Convergence Under Quiescent Streams]
\label{cor:quiescent}
If the stream $\mathcal{S}$ is finite, then there exists a valid state
$\sigma^*$ such that both processors converge to~$\sigma^*$: there exists $T$
with $\Psi_{P_1}(t) = \Psi_{P_2}(t) = \sigma^*$ for all $t \geq T$.
\end{corollary}
\begin{proof}
By eventual delivery, there exists $T_N$ such that
$E_{P_1}(t) = E_{P_2}(t) = \mathcal{S}$ for all $t \geq T_N$.  By
\Cref{thm:convergence}(c), both processors are in the same state.  Since no
new events arrive, this state is fixed from~$T_N$ onward.
\end{proof}

\begin{remark}[Infinite Streams]
\label{rem:infinite-streams}
For infinite streams, the received event sets grow monotonically
toward~$\mathcal{S}$ but may differ at any given time.  When event sets
coincide, \Cref{thm:convergence}(c) guarantees agreement.  Eventual delivery
ensures that any event causing disagreement is transient.  The agreed-upon
state may continue to change as new events arrive---the guarantee is perpetual
agreement, not convergence to a fixed point.  \Cref{cor:quiescent} recovers
fixed-point convergence when the stream is finite.
\end{remark}

%========================================
\section{Necessity of the Structural Conditions}
\label{sec:necessity}

We show that both UBC and CC are necessary for their respective guarantees:
UBC for uniform complexity, CC for agreement.

\subsection{Necessity of UBC}

UBC is not merely convenient; it is necessary for a uniform bound on
compensation depth, and thus for the constant per-event compensation overhead
used in the complexity analysis (\Cref{sec:complexity}).  Without it,
compensation depth can grow without limit as event sets grow.

\begin{theorem}[Unbounded Compensation Depth Without UBC]
\label{thm:necessity}
For any $M \in \mathbb{N}$, there exists a registry $R_\infty$ and a finite
event set~$E$ such that some reduction sequence from $(\sigma_0, E)$ in the
governance rewrite system contains more than $M$ compensation steps.
Consequently, no uniform bound on compensation depth exists, and UBC fails.
\end{theorem}

\begin{proof}
Define $R_\infty$ with $\Sigma_{R_\infty} = \mathbb{Z}$, the single invariant
$\psi(\sigma) = \top \iff \sigma = 0$, initial state $\sigma_0 = 0$, and
compensation $\rho(\sigma) = \sigma + 1$ for $\sigma < 0$ and
$\rho(\sigma) = \sigma - 1$ for $\sigma > 0$ (so compensation moves
toward~$0$).

For each $n \in \mathbb{N}$, consider the single-event set $E_n = \{e_n\}$
where $\mathrm{apply}(e_n, 0) = -n$.  Under the processor model
(\Cref{def:processor}), the processor applies $e_n$ to reach state $-n$, then
compensates incrementally: $-n \to -(n-1) \to \cdots \to -1 \to 0$.  This
requires exactly $n$ compensation steps.

For any proposed bound~$M$, choose $n = M + 1$.  The compensation episode
after this single event contains $n > M$ steps.  Therefore no finite~$M$
uniformly bounds compensation depth, and UBC fails.

The measure $\Phi(\sigma) = |\sigma|$ strictly decreases under~$\rho$, so WFC
holds and compensation terminates for each fixed~$E_n$.  The pathology is not
non-termination but \emph{unbounded} termination: the compensation depth
required to restore validity grows without limit.
\end{proof}

\begin{remark}[What Fails]
The counterexample satisfies WFC (compensation terminates for every
fixed~$E_n$) but violates UBC (compensation depth is~$n$, unbounded).  CC also
holds: there is only one event per set, so order independence is vacuous, and
compensation absorption holds because $\rho$ moves deterministically
toward~$0$.  The pathology is purely one of unbounded depth, which invalidates
the uniform step bound and the constant per-event compensation cost in
\Cref{thm:complexity}.  Finite semantic state spaces prevent this by
bounding~$\Phi$ (\Cref{lem:finite-implies-ubc}).
\end{remark}

\subsection{Necessity of CC}

CC is necessary for agreement: without it, processors applying the same events
in different causality-respecting orders can reach different valid states.

\begin{proposition}[Disagreement Without CC]
\label{prop:cc-necessary}
There exists a registry satisfying WFC and UBC, and a finite event set~$E$
with two causality-respecting orderings, such that the governance rewrite
system reaches different valid normal forms under the two orderings.
\end{proposition}

\begin{proof}
Let $\Sigma_R = \{A, B, C, D\}$ with valid states $\{A, B\}$ and
compensation $\rho_R(C) = A$, $\rho_R(D) = B$.  UBC holds with $M = 1$.

Define two events $e_1, e_2$ with $\mathrm{deps}(e_1) = \mathrm{deps}(e_2) =
\varnothing$, so both are enabled from the initial configuration
$(A, \{e_1, e_2\})$ and are causally independent under
\Cref{def:causal-order}.  Their application is:
\[
\begin{array}{c|cccc}
  & A & B & C & D \\ \hline
e_1 & C & C & C & D \\
e_2 & D & C & C & D
\end{array}
\]

Evaluate CC1 at $\sigma = A$.  Under the processor model
(\Cref{def:processor}), each processor normalizes after each application,
so the comparison is between $\rho_R^*(A_{e_2}(\rho_R^*(A_{e_1}(A))))$ and
the symmetric order.

\emph{Path 1} ($e_1$ then $e_2$): $\mathrm{apply}(e_1, A) = C$,
$\rho_R^*(C) = A$, $\mathrm{apply}(e_2, A) = D$, $\rho_R^*(D) = B$.
Normal form: $B$.

\emph{Path 2} ($e_2$ then $e_1$): $\mathrm{apply}(e_2, A) = D$,
$\rho_R^*(D) = B$, $\mathrm{apply}(e_1, B) = C$, $\rho_R^*(C) = A$.
Normal form: $A$.

Since $A \neq B$, CC1 is violated.  Two processors consuming
$\{e_1, e_2\}$ in different causality-respecting orders reach different valid
states despite WFC and UBC holding.
\end{proof}

\begin{remark}
The key feature is that $\rho_R$ maps distinct invalid states to
\emph{distinct} valid states ($C \mapsto A$, $D \mapsto B$), creating
path-dependent repair.  Contrast this with the order fulfillment example
(\Cref{sec:example}), where compensation preserves approval intent via the
held state, ensuring path-independent repair.  CC captures precisely the
boundary between these cases.
\end{remark}

\begin{remark}
The UBC counterexample (\Cref{thm:necessity}) satisfies CC but violates UBC;
the CC counterexample (\Cref{prop:cc-necessary}) satisfies WFC and UBC but
violates CC.  Together, they show that each condition is independently
necessary for its respective guarantee.
\end{remark}

%========================================
\section{Complexity Analysis}
\label{sec:complexity}

\begin{theorem}[Convergence Complexity]
\label{thm:complexity}
For a stream prefix of length $n$ under a registry satisfying UBC with
bound~$M$ and $k = |I_R|$ invariants, the model-level cost of
registry-governed stream processing is $O(n)$ for event application,
validation, and compensation combined.  If causality maintenance is
implemented with a priority queue, the total cost is $O(n \log n)$, where $k$
and $M$ are treated as constants of system design.
\end{theorem}

\begin{proof}
We analyze the per-event cost, separating model-level costs (which follow from
the formalism) from scheduling costs (which depend on implementation).

\textbf{Model-level per-event costs.}  For each event $e_i$
($1 \leq i \leq n$):
\begin{itemize}[leftmargin=*,nosep]
\item \emph{Event application}: $O(1)$ (state update).
\item \emph{Invariant validation}: $O(k) = O(1)$ (evaluate each~$\psi_j$).
\item \emph{Compensation}: At most $M$ steps, each $O(1)$.  Since $M$ is
  bounded by UBC, this is $O(1)$.
\end{itemize}
Registry governance therefore adds $O(1)$ overhead per event beyond what the
base stream processor already performs.

\textbf{Scheduling overhead.}  Maintaining a set of enabled events and
resolving causal dependencies requires implementation-dependent data
structures.  With a priority queue keyed by timestamp and identifier, each
event insertion and dependency-resolution check costs $O(\log i)$, yielding:

\textbf{Total cost.}
\[
T(n) = \sum_{i=1}^n \bigl(\underbrace{O(1)}_{\text{model}} +
  \underbrace{O(\log i)}_{\text{scheduling}}\bigr)
     = O(n) + O\!\left(\sum_{i=1}^n \log i\right)
     = O(n \log n).
\]

The dominant term is causality maintenance.  Validation and compensation
contribute no asymptotic overhead.  Systems with native causal ordering (e.g.,
single-partition Kafka topics) eliminate the $O(\log i)$ term, reducing the
bound to~$O(n)$.
\end{proof}

%========================================
\section{Federated Convergence}
\label{sec:federation}

The convergence theorem (\Cref{thm:convergence}) applies within a single
registry.  The product composition theorem (\Cref{thm:product}) extends it
to independent registries with disjoint invariants.  But real multi-registry
systems have \emph{dependencies} across organizational boundaries: a
manufacturer's specifications constrain supplier capabilities; a regulator's
rules constrain bank behavior.  We now extend convergence to registries
connected by morphisms encoding such constraints.

The central observation is that a federation of registries connected by
morphisms \emph{is} a registry in the sense of \Cref{sec:model}: the
federated state space is the product of component state spaces, morphism
consistency conditions are additional invariants, and morphism repair is
additional compensation.  Federated convergence then follows from the
single-registry theory if the federated registry satisfies WFC and CC.  WFC,
and with it a unique repair normal form, is derived from the network's
structure (\Cref{lem:fed-termination}).  CC is not: a target event can read a
shared value that morphism repair overwrites, and then the components' own CC
does not carry over to the federation (\Cref{rem:why-c1c2}).  We isolate the
two per-edge conditions that are needed, C1 and C2 (\Cref{def:c1c2}), prove
that they give federated CC and convergence (\Cref{thm:fed-cc,thm:fed-convergence}),
and prove that their reachable forms are exact.

\subsection{Registry Morphisms}

\begin{definition}[Registry Morphism]
\label{def:morphism}
Given registries $R_A$ and $R_B$, a \emph{registry morphism}
$\phi: R_A \to R_B$ consists of a \emph{decomposition} of $B$'s state
space---an isomorphism
$\langle \pi_B^{\mathrm{sh}}, \pi_B^{\mathrm{loc}} \rangle:
\Sigma_B \xrightarrow{\;\sim\;} S_B \times L_B$
with reconstruction $\mathrm{merge}_B: S_B \times L_B \to \Sigma_B$
satisfying $\mathrm{merge}_B(\pi_B^{\mathrm{sh}}(\sigma),\,
\pi_B^{\mathrm{loc}}(\sigma)) = \sigma$---and a function
$\phi: \Sigma_A \to S_B$ mapping source states to \emph{shared} target
states.  Write
$\sigma_B[\mathrm{sh} \mapsto s] \defeq
\mathrm{merge}_B(s,\, \pi_B^{\mathrm{loc}}(\sigma_B))$.

The morphism satisfies:
\begin{enumerate}[leftmargin=*,nosep,label=(M\arabic*)]
\item \textbf{Validity preservation under overwrite}: If
  $V_A(\sigma_A) = \top$ and $V_B(\sigma_B) = \top$, then
  $V_B(\sigma_B[\mathrm{sh} \mapsto \phi(\sigma_A)]) = \top$.
\end{enumerate}
When $\pi_B^{\mathrm{sh}} = \mathrm{id}$ ($L_B$ trivial), the morphism
constrains $B$ completely (\emph{total morphism}).
\end{definition}

\begin{remark}
M1 is not standard function-level validity preservation; it is validity
preservation under shared-component overwrite.  A sufficient condition: if
$V_B$ factors as
$V_B(\sigma) \iff V_B^{\mathrm{sh}}(\pi_B^{\mathrm{sh}}(\sigma)) \wedge
V_B^{\mathrm{loc}}(\pi_B^{\mathrm{loc}}(\sigma))$ (shared and local validity
independent), then M1 reduces to
$V_A(\sigma_A) = \top \Rightarrow
V_B^{\mathrm{sh}}(\phi(\sigma_A)) = \top$.
The shared/local decomposition models partial authority: a manufacturer
constrains product specifications (shared) but not supplier pricing (local).

M1 is a nonvacuous constraint: its antecedent requires both a valid source
and a valid target state, and WFC (\Cref{ax:wfc}) guarantees every component
registry has a nonempty valid set, since iterated compensation reaches a valid
state from any state. A registry whose invariants admit no valid state is
therefore excluded at the component level, not silently satisfying M1 here.
\end{remark}

\begin{definition}[Tree-Shaped Morphism Network]
\label{def:network}
A \emph{morphism network} $\mathcal{N} = (\{R_i\}_{i \in \mathcal{I}},
\{\phi_{ij}\}_{(i,j) \in E})$ is \emph{tree-shaped} if $(\mathcal{I}, E)$ is
a directed forest: each connected component is a tree rooted at a
\emph{source} registry (no incoming morphisms), and each non-source has
exactly one incoming edge.
\end{definition}

\subsection{Federated Construction}

\begin{definition}[Federated Registry]
\label{def:fed-registry}
Given a tree-shaped morphism network $\mathcal{N}$, the federated registry
$R_\Fed$ has state space $\Sigma_\Fed = \prod_i \Sigma_{R_i}$ with
projections $\pi_i(\bar{\sigma}) = \sigma_i$.  Invariants comprise:
(i)~\emph{local invariants}---each $\psi_j \in I_{R_i}$ lifted to the product
as $\hat{\psi}_{i,j}(\bar{\sigma}) = \psi_j(\sigma_i)$---and
(ii)~\emph{morphism invariants}---for each $(i,j) \in E$,
$\mu_{ij}(\bar{\sigma}) = \top \iff
\phi_{ij}(\sigma_i) = \pi_j^{\mathrm{sh}}(\sigma_j)$.
\end{definition}

\begin{definition}[Federated Normalization Operator]
\label{def:fed-comp}
Let $i_1, \ldots, i_n$ be a topological ordering of $\mathcal{I}$.
The operator $\rho_\Fed: \Sigma_\Fed \to \Sigma_\Fed$ acts in two phases:

\textbf{Phase 1 (Local normalization):} For each~$i$, replace $\sigma_i$ with
$\rho_{R_i}^*(\sigma_i)$.

\textbf{Phase 2 (Morphism repair):} Process edges in topological order.
For $(i,j) \in E$, if $\phi_{ij}(\sigma_i) \neq
\pi_j^{\mathrm{sh}}(\sigma_j)$, set
$\sigma_j \mapsto \sigma_j[\mathrm{sh} \mapsto \phi_{ij}(\sigma_i)]$.

Local normalization precedes morphism repair because M1 guarantees that
overwriting the shared component of a locally valid target with the morphism
image of a locally valid source preserves target validity.
\end{definition}

\begin{definition}[Federated Governance Rewrite System]
\label{def:fed-grs}
A \emph{federated event} $e$ targets one registry $R_{m(e)}$: applying it
replaces $\sigma_{m(e)}$ by $\apply(e, \sigma_{m(e)})$ and leaves every other
component unchanged, so cross-registry effects arise only through morphism
violation and repair.  Write $V_\Fed$ for the conjunction of all local and
morphism invariants (\emph{federal validity}).  The \emph{federated governance
rewrite system} $\mathcal{G}_\Fed$ operates on configurations
$(\bar{\sigma}, B)$ as in \Cref{def:grs}, with two rules:
\begin{enumerate}[leftmargin=*,nosep]
\item \textbf{Compensation step:} $(\bar{\sigma}, B) \to (\rho_\Fed(\bar{\sigma}), B)$
  whenever $V_\Fed(\bar{\sigma}) = \bot$.
\item \textbf{Guarded apply step:}
  $(\bar{\sigma}, B) \to (\apply(e, \bar{\sigma}), B \setminus \{e\})$ whenever
  $e \in B$ is enabled and $V_\Fed(\bar{\sigma}) = \top$.
\end{enumerate}
The guard says that events are applied to normalized states: a federated
processor compensates, then applies the next event.  An apply step followed by
the compensation it triggers is the \emph{governed step}
$\bar{\sigma} \mapsto \rho_\Fed(\apply(e, \bar{\sigma}))$.  The unguarded
system, in which apply steps may also fire at federally invalid states, is
discussed in \Cref{rem:unguarded}.
\end{definition}

\paragraph{Notation.}
For an event $e$ of registry~$j$ write
$\mathrm{sig}_e \defeq \rho_{R_j}^* \circ \apply(e, -)$ (the component step: the
event, then component compensation).  For a registry $j$ and a family
$z = (z_i)_{i \in S(j)}$ of states of its sources write
$\mathrm{ow}_z(x) \defeq x[\mathrm{sh} \mapsto \phi_{ij}(z_i)]$ (the morphism
repair of~$j$ against~$z$; for a multi-source target the resolved value
$\Gamma_j(z)$ of \Cref{sec:multisource} replaces $\phi_{ij}(z_i)$; for a source
registry $\mathrm{ow}_z$ is the identity).  A state $b$ of~$j$ is
\emph{consistent with}~$z$ if $b = \mathrm{ow}_z(b)$.  Two federated events are
\emph{federated-independent} if they target different registries, or target the
same registry and are causally independent.

\begin{definition}[Cross-Registry Compensation Commutativity]
\label{def:c1c2}
A morphism network satisfies:
\begin{enumerate}[leftmargin=*,nosep]
\item[\textbf{(C1)}] \textbf{Cross-registry CC.}\; For every event $e$ of
  registry~$j$, all valid source families $z, z'$ of~$j$, and every valid
  state $b$ of~$j$ consistent with~$z'$:
  \[
  \mathrm{ow}_z(\mathrm{sig}_e(\mathrm{ow}_z(b))) \;=\; \mathrm{ow}_z(\mathrm{sig}_e(b)).
  \]
\item[\textbf{(C2)}] \textbf{Repaired CC.}\; For all causally independent events
  $e_1, e_2$ of the same registry~$j$, every valid source family~$z$ of~$j$, and
  every valid $b$ consistent with~$z$:
  \[
  \mathrm{ow}_z(\mathrm{sig}_{e_2}(\mathrm{ow}_z(\mathrm{sig}_{e_1}(b))))
  \;=\;
  \mathrm{ow}_z(\mathrm{sig}_{e_1}(\mathrm{ow}_z(\mathrm{sig}_{e_2}(b)))).
  \]
\end{enumerate}
\end{definition}

C1 says that the outcome of a target event, after repair against the current
sources~$z$, does not depend on whether the sources moved from $z'$ to $z$
before or after the event.  The event may read and write the shared component;
the final overwrite appears on both sides.  C2 is CC for the step ``event, then
repair''.  For a source registry ($\mathrm{ow}$ the identity) C1 is trivial and
C2 is exactly the registry's own CC1 at valid states.  Both quantify over one
edge (a target and its sources), so on finite state spaces both are checked by
enumeration, like M1.

\begin{remark}[Why component CC is not enough]
\label{rem:why-c1c2}
For a target, neither C1 nor C2 follows from the target's own CC and M1.  In a
two-registry tree satisfying M1 and component WFC and CC, a target event that
copies a shared flag (which the morphism overwrites) into a local counter ends
with the counter at~$0$ or at~$1$ depending on whether a source event is
propagated first: C1 fails and two event orders diverge (Coq:
\texttt{audit\_counterexample}, and against this paper's exact hypotheses
\texttt{fed\_thm\_fed\_convergence\_refuted}).  A second tree satisfies C1 and
component CC (the target's two events commute locally) but violates C2,
because the repair between the two events erases what the first wrote into a
shared variable; it diverges as well (\texttt{c2\_counterexample},
\texttt{fed\_c2\_paper\_counterexample}).  So neither condition can be dropped.
\end{remark}

\subsection{The Authority Argument}

In a directed morphism $\phi: R_A \to R_B$,
registry~$A$ is \emph{authoritative} over $B$'s shared component: in a
federally valid state, $B$'s shared state is completely determined by $A$'s
state.  Throughout, \textbf{(H)} denotes the standing hypotheses: a tree-shaped
morphism network (\Cref{def:network}) whose morphisms satisfy M1 and whose
component registries satisfy WFC and CC.

\begin{lemma}[Authority Determination]
\label{lem:authority}
Assume (H).
\begin{enumerate}[leftmargin=*,nosep,label=(\alph*)]
\item If $(\bar{\sigma}^0, B) \to^* (\bar{\nu}, \varnothing)$ in
  $\mathcal{G}_\Fed$ with $\bar{\nu}$ federally valid, then for every source
  registry~$i$, $\nu_i$ is the unique normal form of the single-registry
  rewrite system of $R_i$ from $(\sigma^0_i, B|_i)$, where $B|_i$ is the set of
  events of $B$ targeting~$i$ (events targeting other registries and morphism
  repair do not modify source components).
\item In every federally valid $\bar{\sigma}$, the shared component of a
  non-source~$j$ with incoming edge from~$i$ satisfies
  $\pi_j^{\mathrm{sh}}(\sigma_j) = \phi_{ij}(\sigma_i)$: it is a function of
  the source state alone.
\item Morphism repair never modifies $\pi_j^{\mathrm{loc}}$.
\end{enumerate}
\end{lemma}

The local component of a target is not determined by component CC: two event
orders can leave it different (\Cref{rem:why-c1c2}).  Under C1 and C2 it is
determined by \Cref{thm:fed-convergence} and computed by \Cref{cor:fed-nf}.
(Coq: (a) \texttt{fed\_lem\_authority\_a}, with instance
\texttt{fed\_lem\_authority\_a\_instance}; (b) \texttt{fed\_lem\_authority\_b};
(c) \texttt{fed\_lem\_authority\_c\_local}; the version~1 convergence clause is
refuted by \texttt{fed\_lem\_authority\_c\_refuted}.  Part~(a) is proved for
the unguarded system and holds for $\mathcal{G}_\Fed$, whose reductions are
among the unguarded ones.)

\begin{lemma}[Single-Round Termination]
\label{lem:fed-termination}
Assume (H).  For every federated state $\bar{\sigma}$, not only locally valid
ones, $\rho_\Fed(\bar{\sigma})$ is federally valid.  Phase~1 on registry~$j$
takes at most $\Phi_j(\sigma_j)$ compensation steps (at most $M_j$ under UBC);
Phase~2 processes each edge once (unique target: each registry's shared
component is set once and not revisited).  Morphism repair cannot reintroduce
local invalidity: M1 guarantees that overwriting the shared component of a
locally valid target with the morphism image of a locally valid source
preserves target validity.  Since Phase~1 establishes local validity before
Phase~2 begins, and Phase~2 only modifies shared components, every
intermediate state of Phase~2 is locally valid.  Moreover $\rho_\Fed$ is the
identity on federally valid states, so it is an idempotent retraction onto
them, and its result does not depend on the topological order chosen.
\end{lemma}

(Coq: \texttt{fed\_lem\_fed\_termination}, \texttt{fed\_phase1\_bound},
\texttt{single\_round}, \texttt{phase2\_invariant}, \texttt{grho\_valid},
\texttt{grho\_id}; order independence and the retraction on the categorical
paper's consistent set: \texttt{cat\_thm\_one\_idempotent},
\texttt{cat\_thm\_one\_image}, \texttt{cat\_thm\_one\_order\_independent}.)
Repair confluence therefore composes with no condition beyond M1: the
\emph{repair normal form} $\rho_\Fed(\bar{\sigma})$ of any state is unique.
Event confluence is the subject of the next two theorems.

\begin{theorem}[Federated CC]
\label{thm:fed-cc}
Assume (H), C1 and C2.  Then the governed step satisfies CC1 at every
federally valid state: for every federally valid $\bar{\sigma}$ and all
federated-independent events $e_1, e_2$,
\[
\rho_\Fed(\apply(e_2, \rho_\Fed(\apply(e_1, \bar{\sigma}))))
\;=\;
\rho_\Fed(\apply(e_1, \rho_\Fed(\apply(e_2, \bar{\sigma})))).
\]
\end{theorem}

\begin{proof}
Let $\bar{\sigma}$ be federally valid and write $\mathrm{gov}_e$ for the
governed step.  On a federally valid state, $\mathrm{gov}_e$ replaces the
component $m(e)$ by $\mathrm{sig}_e$ of it and then re-repairs every target in
topological order against its finalized sources; by \Cref{lem:authority}(b)
every shared component is the image of its sources.  Process the registries in
topological order and compare the two sides component by component.  A
registry that is neither targeted by $e_1$ or $e_2$ nor downstream of one is
unchanged on both sides.  If $e_1, e_2$ target different registries
$m_1 \neq m_2$, the component $m_1$ is, on one side, $\mathrm{ow}_z$ applied to
$\mathrm{sig}_{e_1}$ of its consistent state and, on the other, the same after
its sources moved from $z'$ to $z$ through $e_2$; C1 at $(e_1, z, z', b)$ says
the two agree, and symmetrically for~$m_2$.  Every other target is the repair of
its own (unchanged) state against sources that agree by induction.  If $e_1,
e_2$ target the same registry, the two sides of C2 at the registry's consistent
state are exactly its two final values, and the downstream registries again
agree by induction.  The formal proof is \texttt{fed\_events\_commute}, and
against (H) \texttt{fed\_thm\_fed\_cc\_corrected\_machine}. \qed
\end{proof}

CC2 is not needed for $\mathcal{G}_\Fed$: no apply step fires at a federally
invalid state, so the critical pair ``apply versus compensation'' does not
arise.  Under (H), C1 and C2 do not give CC2 of the unguarded system
(\Cref{rem:unguarded}).  In the guarded proof the components' own CC enters only
through C2, which for a source registry is that registry's CC1.

\begin{theorem}[Federated Normalization Confluence]
\label{thm:fed-convergence}
Assume (H) and let $\bar{\sigma}^0$ be federally valid.
\begin{enumerate}[leftmargin=*,nosep,label=(\alph*)]
\item If C1 and C2 hold, every configuration $(\bar{\sigma}^0, B)$ has a unique
  normal form in $\mathcal{G}_\Fed$, and it is federally valid: any two
  processors that consume the same federated events from $\bar{\sigma}^0$,
  applying each event to a normalized state, converge to the same federally
  valid state.
\item \textbf{(Exact condition.)}  Every configuration $(\bar{\sigma}^0, B)$ has
  a unique normal form in $\mathcal{G}_\Fed$ if and only if C1 holds at every
  \emph{reachable witness} $(e, z, z', b) = (e,\, \mathrm{gov}_a(\bar{\sigma}),\,
  \bar{\sigma},\, \sigma_{m(e)})$ and C2 at every reachable witness
  $(e_1, e_2, z, b) = (e_1, e_2, \bar{\sigma}, \sigma_{m(e_1)})$, where
  $\bar{\sigma}$ ranges over the states reachable from $\bar{\sigma}^0$ by
  governed steps and $a$ over events of registries other than $m(e)$.
  Equivalently, federated-independent events commute after re-normalization at
  every reachable state (the global condition GC).
\end{enumerate}
\end{theorem}

\begin{proof}
\emph{Termination.}  By \Cref{lem:fed-termination}, $\rho_\Fed$ maps any
federated state to a federally valid one in a single application, so the
federated measure $\Phi_\Fed(\bar{\sigma}) = \mathbf{1}[V_\Fed(\bar{\sigma}) =
\bot] \in \{0, 1\}$ satisfies WFC, and the lexicographic order $(|B|,
\Phi_\Fed)$ gives termination as in \Cref{lem:termination}.

\emph{(a).}  Local confluence of $\mathcal{G}_\Fed$: two compensation steps
coincide ($\rho_\Fed$ is deterministic); an apply step and a compensation step
are never enabled together (the guard); two apply steps at a federally valid
state close after one compensation each by \Cref{thm:fed-cc}.  Newman's Lemma
gives unique normal forms, which are federally valid by the definition of
normal form.

\emph{(b).}  Every reachable state of $\mathcal{G}_\Fed$ that is federally valid
is the result of a sequence of governed steps from $\bar{\sigma}^0$.  At such a
state the two sides of C1 at the reachable witness are the registry $m(e)$ after
the governed runs $e;a$ and $a;e$ (which differ by one allowed swap), and the
two sides of C2 are the registry after $e_1;e_2$ and $e_2;e_1$.  So a failing
witness is an observable divergence, and conversely, if every reachable witness
passes, independent events commute at every reachable state and the argument of
(a) applies with ``reachable'' in place of ``every''. \qed
\end{proof}

(Coq: (a) \texttt{fed\_thm\_fed\_convergence\_guarded}, in the machine model
\texttt{fed\_interleavings\_converge} and \texttt{fed\_permutations\_converge};
(b) \texttt{fed\_thm\_fed\_convergence\_exact}, \texttt{fed\_guarded\_exact},
\texttt{fed\_exact}; with any number of intervening events in place of one,
\texttt{fed\_exact\_full}; GC: \texttt{acyclic\_gc\_iff},
\texttt{gc\_iff\_reach}; static C1 and C2 imply the reachable forms:
\texttt{static\_c1\_c2\_gc}; the two divergence constructions:
\texttt{conv\_c1\_diverge}, \texttt{conv\_c2\_diverge}.  Non-vacuity: a supply
federation satisfying (H), C1 and C2, \texttt{fed\_supply\_paper\_instance}.)
Coq's buffers carry no causal dependencies (every buffered event is enabled);
with dependencies the reductions are a subset, so~(a) carries over.

\begin{remark}[Static and reachable checks]
\label{rem:static-reachable}
Static C1 and C2 (\Cref{def:c1c2}) quantify over all valid source families and
all valid consistent target states; they are per-edge and are what a build-time
verifier enumerates.  Their reachable forms in \Cref{thm:fed-convergence}(b)
are exact but require exploring runs.  The naive converse ``C1 fails, so some
two interleavings diverge'' is false: a federation can violate C1 at a pair of
source states that no run moves between, while every interleaving converges
(\texttt{naive\_converse\_fails}).
\end{remark}

\begin{remark}[The unguarded system]
\label{rem:unguarded}
Version~1 of this paper used the unguarded system, in which an apply step may
fire at a federally invalid state.  There C1 and C2 are not sufficient: a tree
satisfying (H), C1 and C2, on which every permutation of governed steps
converges and $\mathcal{G}_\Fed$ has unique normal forms, nevertheless reaches
two normal forms in the unguarded system from one federally valid start,
because two target events are applied before compensation and the second reads
the shared value the first wrote (\texttt{fed\_grs\_c1\_c2\_insufficient}).
For the unguarded system the exact condition is CC1 and CC2 of the governed
step on the states reachable from the start (\texttt{fed\_grs\_exact}, an
instance of the single-registry exact converse), which implies the reachable
C1/C2 condition but not conversely (\texttt{fed\_grs\_un\_c1\_c2}).  A clean
sufficient condition is \textbf{XU}: C1 for \emph{every} valid target state~$b$,
not only consistent ones.  Under (H) and XU the unguarded system satisfies CC1
and CC2 at every state and has unique normal forms
(\texttt{fed\_thm\_fed\_cc\_corrected},
\texttt{fed\_thm\_fed\_convergence\_corrected}, \texttt{grs\_unique\_nf}).  XU
with component CC implies C1 and C2 (\texttt{xu\_implies\_c1\_c2}), and XU is
also what a distributed model with explicit propagation steps needs
(\texttt{dist\_interleavings\_converge}).
\end{remark}

\begin{corollary}[Constructive Normal Form]
\label{cor:fed-nf}
Assume (H), C1 and C2, and let $\bar{\sigma}$ be federally valid.  For every
event sequence, the federated normal form $\bar{Z}$ is the unique solution,
computed registry by registry in topological order, of the following
equations.  For registry~$j$ with events $e_1, \ldots, e_m$ (in sequence order),
\[
x_0 = \mathrm{ow}_{Z}(\sigma_j), \qquad
x_t = \mathrm{ow}_{Z}(\mathrm{sig}_{e_t}(x_{t-1})), \qquad
Z_j = \mathrm{ow}_{Z}(x_m),
\]
where $\mathrm{ow}_Z$ reads only the finalized sources of~$j$.  So: compute each
source's single-registry run on its own events ($\mathrm{ow}$ is the identity
there), then, in topological order, run each target's events with the repair
against its finalized sources interleaved after \emph{every} event.  The
product state space is never materialized.
\end{corollary}

(Coq: \texttt{fed\_cor\_fed\_nf\_corrected}, \texttt{fed\_nf\_constructive}.)
Under XU the interleaved repairs can be dropped and the recipe of version~1
holds: run the target's events as a single registry, then overwrite its shared
component once (\texttt{fed\_nf\_recipe\_xu}).  Under C1 and C2 alone that
recipe is false (\texttt{fed\_cor\_fed\_nf\_refuted}: it gives a different state
from every event order).

\begin{remark}[Compensation Compatibility is Derived]
One might expect morphisms to satisfy
$\phi(\rho_A^*(\sigma)) = \rho_B^*(\phi(\sigma))$.  The authority argument
shows this is unnecessary: $\rho_\Fed$ normalizes locally first, then projects
via~$\phi$.  The composition ``normalize source, then project'' produces a
valid shared component by M1 alone.  We never need the reverse.
\end{remark}

\begin{remark}[Relationship to Product Composition]
\Cref{thm:product} handles the special case $E = \varnothing$: independent
registries with no morphism constraints.  There C1 is trivial and C2 is each
component's own CC1, so \Cref{thm:fed-convergence} reduces to product
composition.  Federated convergence generalizes it to registries connected by
directed semantic constraints.
\end{remark}

\subsection{Necessity of the Structural Conditions}

\begin{proposition}[Necessity of Acyclicity]
\label{prop:cycle-necessary}
There exists a cyclic morphism network where each component satisfies WFC and
CC, all morphisms satisfy M1, but federated compensation cycles.
\end{proposition}

\begin{proof}
$R_A = R_B$: $\Sigma = \{0, 1\}$, both valid, trivial compensation.
Morphisms: $\phi_{AB}(x) = 1 - x$ (flip), $\phi_{BA} = \mathrm{id}$.  Both
satisfy M1.  From $(0, 0)$: repair on $\phi_{AB}$ gives $(0, 1)$; repair on
$\phi_{BA}$ gives $(1, 1)$; repair on $\phi_{AB}$ gives $(1, 0)$; repair on
$\phi_{BA}$ gives $(0, 0)$.  Cycle.  Every repair step is forced (each state
violates exactly one morphism invariant), no state is federally valid, and
compensation never terminates from any state (Coq:
\texttt{prop\_cycle\_necessary}, \texttt{cycle\_paper\_trace}). \qed
\end{proof}

\begin{proposition}[Necessity of M1]
\label{prop:m1-necessary}
There exists an acyclic network where each component satisfies WFC and CC, but
a morphism violates M1, and federated compensation oscillates.
\end{proposition}

\begin{proof}
$R_A$: $\Sigma_A = \{0\}$, trivial.  $R_B$: $\Sigma_B = \{0, 1\}$,
$V_B(0) = \top$, $V_B(1) = \bot$, $\rho_B(1) = 0$.  Total morphism
$\phi_{AB}(0) = 1$ (violates M1).  From $(0, 0)$: morphism repair sets
$\sigma_B = 1$; local compensation sets $\sigma_B = 0$; morphism repair
sets $\sigma_B = 1$.  Oscillation between morphism repair and local
compensation: no state is federally valid and compensation never terminates
(Coq: \texttt{prop\_m1\_necessary}, \texttt{m1\_paper\_trace}).  In
particular one application of $\rho_\Fed$ from $(0,0)$ returns $(0,1)$, which
is not federally valid, so \Cref{lem:fed-termination} needs M1
(\texttt{m1\_single\_round\_fails}). \qed
\end{proof}

\Cref{rem:why-c1c2} gives the corresponding counterexamples for C1 and C2:
with acyclicity, M1 and component WFC and CC all holding, dropping either check
lets two event orders reach different federally valid states.

\begin{remark}[The Multi-Source Boundary]
\label{rem:multisource}
The tree-shaped restriction excludes networks where a registry has multiple
incoming morphisms from independent sources.  When two morphisms
$\phi_{ij}$ and $\phi_{kj}$ both target registry~$j$, a federally valid
state requires $\phi_{ij}(\sigma_i) = \phi_{kj}(\sigma_k) =
\pi_j^{\mathrm{sh}}(\sigma_j)$---constraining the \emph{sources} to produce
compatible projections.  The authority argument breaks: no single source
determines the target.  Resolving this requires additional machinery
(composition consistency along common ancestors, target agreement predicates,
or conflict resolution operators).  We develop the last of these---conflict
resolution operators---in \Cref{sec:multisource}, recovering the unique repair
normal form on any acyclic network and, under C1 and C2, event convergence.
\end{remark}

\subsection{Multi-Source Resolution}
\label{sec:multisource}

\Cref{rem:multisource} identified the obstruction to networks in which a
registry has several incoming morphisms: no single source determines the
target's shared component.  We remove the tree restriction---admitting any
acyclic network---by replacing the per-edge authority function at a multi-source
target with a single \emph{resolution operator}.  The authority argument then
generalizes: a target's shared component becomes a deterministic function not of
one source but of all of them.

\begin{definition}[Resolution Operator]
\label{def:resolver}
Let $j$ be a registry with source set $S(j) = \{i : (i,j) \in E\}$.  A
\emph{resolution operator} for~$j$ is a map
\[
  \Gamma_j : \textstyle\prod_{i \in S(j)} \Sigma_{R_i} \longrightarrow S_j
\]
from the source states to $j$'s shared component, satisfying:
\begin{enumerate}[leftmargin=*,nosep,label=(R\arabic*)]
\item \textbf{Source determinacy}: $\Gamma_j$ is a function of the source states
  alone; it does not read $\pi_j^{\mathrm{loc}}(\sigma_j)$.
\item \textbf{Validity preservation}: for all $(\sigma_i)_{i \in S(j)}$ with
  each $V_{R_i}(\sigma_i) = \top$ and all $\sigma_j$ with
  $V_{R_j}(\sigma_j) = \top$,
  \[
    V_{R_j}\!\left(\sigma_j[\mathrm{sh} \mapsto
      \Gamma_j((\sigma_i)_{i \in S(j)})]\right) = \top .
  \]
\end{enumerate}
When $|S(j)| = 1$, taking $\Gamma_j(\sigma_i) = \phi_{ij}(\sigma_i)$ recovers the
authority morphism and (R2) is exactly M1.  A resolution operator thus strictly
generalizes the single-source morphism image: it expresses merges---priority,
conjunction, most-restrictive-wins---that no single authority can.
As with M1, (R2) is nonvacuous: WFC (\Cref{ax:wfc}) on each source and on~$j$
makes both antecedents satisfiable.
\end{definition}

\begin{definition}[Resolved Network]
\label{def:resolved-network}
A \emph{resolved network} is an acyclic morphism network in which every target
with $|S(j)| > 1$ carries a resolution operator $\Gamma_j$; single-source targets
retain their authority morphism.  The federated registry is as in
\Cref{def:fed-registry}, with the morphism invariants for a multi-source target
replaced by the \emph{resolution invariant}
\[
  \mu_j(\bar{\sigma}) = \top \iff
  \pi_j^{\mathrm{sh}}(\sigma_j) = \Gamma_j((\sigma_i)_{i \in S(j)}) .
\]
\end{definition}

The normalization operator $\rho_\Fed$ carries over from \Cref{def:fed-comp}
unchanged, with Phase~2 setting each multi-source target's shared component in a
single step: for target~$j$ in topological order, if
$\pi_j^{\mathrm{sh}}(\sigma_j) \neq \Gamma_j((\sigma_i)_{i \in S(j)})$, set
$\sigma_j \mapsto \sigma_j[\mathrm{sh} \mapsto \Gamma_j((\sigma_i)_{i \in S(j)})]$.
Acyclicity guarantees a topological ordering in which every source
$i \in S(j)$ precedes~$j$.

\begin{lemma}[Single-Round Termination, Resolved]
\label{lem:resolved-termination}
On a resolved network whose single-source morphisms satisfy M1, whose
resolution operators satisfy (R1)--(R2), and whose components satisfy WFC and
CC, a single application of $\rho_\Fed$ to any federated state produces a
federally valid state.
\end{lemma}

\begin{proof}
Identical to \Cref{lem:fed-termination}, with the single-source repair
$\phi_{ij}(\sigma_i)$ replaced by $\Gamma_j((\sigma_i)_{i \in S(j)})$ and M1 by (R2).
Phase~1 establishes local validity; Phase~2 sets each target's shared component
exactly once (its resolution operator consumes all incoming edges in one
evaluation, and sources are finalized first by acyclicity).  By (R2),
overwriting a locally valid target's shared component with the resolved value of
locally valid sources preserves target validity, so Phase~2 never reintroduces
invalidity.  (Coq: \texttt{fed\_lem\_resolved\_termination}.) \qed
\end{proof}

The rewrite system $\mathcal{G}_\Fed$ (\Cref{def:fed-grs}), the notation
$\mathrm{ow}_z$ (with $\Gamma_j(z)$ as the resolved value) and the checks C1 and
C2 (\Cref{def:c1c2}) carry over unchanged.  Write \textbf{(H$_R$)} for the
hypotheses of \Cref{lem:resolved-termination}: an acyclic resolved network,
M1 on single-source morphisms, (R1)--(R2) on resolvers, and WFC and CC on every
component.  \Cref{lem:authority} holds under (H$_R$) with
$\Gamma_j((\sigma_i)_{i \in S(j)})$ in place of $\phi_{ij}(\sigma_i)$ in~(b)
(the Coq statements are proved for resolved networks, the tree case being an
instance).

\begin{theorem}[Federated Convergence with Resolution]
\label{thm:resolved-convergence}
Assume (H$_R$) and let $\bar{\sigma}^0$ be federally valid.
\begin{enumerate}[leftmargin=*,nosep,label=(\alph*)]
\item If C1 and C2 hold, every configuration $(\bar{\sigma}^0, B)$ has a unique
  normal form in $\mathcal{G}_\Fed$, and it is federally valid: any two
  processors consuming the same federated events from $\bar{\sigma}^0$ converge
  to the same federally valid state.
\item Every configuration $(\bar{\sigma}^0, B)$ has a unique normal form in
  $\mathcal{G}_\Fed$ if and only if C1 and C2 hold at the reachable witnesses
  of \Cref{thm:fed-convergence}(b).
\end{enumerate}
\end{theorem}

\begin{proof}
As for \Cref{thm:fed-convergence}.  \Cref{lem:resolved-termination} gives WFC
with the trivial measure $\Phi_\Fed \in \{0,1\}$.  The proof of
\Cref{thm:fed-cc} uses only that every target is repaired once, in topological
order, against its finalized sources, by a repair that reads the sources alone
(R1) and preserves validity (R2); it applies verbatim with $\Gamma_j$ in place
of $\phi_{ij}$.  Newman's Lemma gives~(a), and the witness construction gives~(b).
(Coq: \texttt{fed\_thm\_resolved\_convergence\_guarded},
\texttt{fed\_thm\_resolved\_convergence\_exact}; for the unguarded system under
XU, \texttt{fed\_thm\_resolved\_convergence\_corrected}; the version~1
statement is refuted by \texttt{fed\_thm\_resolved\_convergence\_refuted}.) \qed
\end{proof}

\begin{corollary}[Constructive Resolved Normal Form]
\label{cor:resolved-nf}
\Cref{cor:fed-nf} extends to (H$_R$) with C1 and C2: compute each source's
single-registry run, then, in topological order, run each target's events with
the repair against its finalized sources (its resolution operator, or authority
morphism) interleaved after every event.  The federated normal form is the
unique solution of these equations.  As in the tree case, the product state
space is never materialized.  (Coq: \texttt{fed\_cor\_resolved\_nf\_corrected}.)
\end{corollary}

\begin{remark}[Verified, not assumed]
Unlike the authority morphism, a resolution operator encodes domain policy:
there is no canonical merge.  The theorem is therefore conditioned on
(R1)--(R2), exactly as the single-source case is conditioned on M1, and on C1
and C2.  All four are finite conditions on one target and its sources: on
finite state spaces they are checked by enumerating
$\prod_{i \in S(j)} \Sigma_{R_i} \times \Sigma_{R_j}$ (with the target's events
for C1 and C2), so a build-time verifier certifies a resolved network as it
certifies WFC and CC for a single registry.  \Cref{thm:fed-convergence} is the
special case in which every target has a single source whose resolution
operator is the morphism image.
\end{remark}

\begin{remark}[Necessity of the resolution conditions]
Acyclicity remains necessary (\Cref{prop:cycle-necessary} is unchanged;
resolution operators do not tame cycles).  (R2) is necessary for the reason M1
is (\Cref{prop:m1-necessary}): a two-source resolver satisfying (R1) that maps
a valid source combination to a shared component violating $V_{R_j}$ makes
resolution repair and local compensation oscillate forever, with no federally
valid state (\texttt{r2\_necessary}).  (R1) is necessary for a unique repair
normal form: a resolver that reads the target's local component, on an acyclic
network where (R2) and component WFC hold, admits two federally valid normal
forms with different shared components from one state, depending on whether
local compensation or resolution repair fires first (\texttt{r1\_necessary}).
\end{remark}

\subsection{Monotone Cycles}
\label{sec:monotone}

Acyclicity (\Cref{def:network,def:resolved-network}) is required because the
necessity counterexamples turn on \emph{non-monotone} repair: the cyclic
counterexample (\Cref{prop:cycle-necessary}) uses the negation morphism
$\phi(x) = 1 - x$, and the M1 counterexample (\Cref{prop:m1-necessary}) an
antitone flip.  We show that when repair is \emph{monotone} on an ordered shared
domain, acyclicity is unnecessary for a unique repair normal form: it is a least
fixed point (Knaster--Tarski), and order-independence of the repair is the
classical convergence of \emph{chaotic iteration}.  Event convergence on a
cycle needs a separate condition, which we also characterize exactly.

\begin{definition}[Ordered Shared Domains]
\label{def:lattice-shared}
A network has \emph{ordered shared domains} if each shared space $S_j$ carries a
partial order $\sqsubseteq_j$ making $(S_j, \sqsubseteq_j)$ a complete lattice
(least element $\bot_j$; all joins exist).  For fixed local components (the
components' Phase~1 normal forms), define the \emph{federated repair
operator} $\Phi : \prod_j S_j \to \prod_j S_j$ by
\[
  \Phi(\bar{s})_j =
  \begin{cases}
    s_j
      & \text{$j$ a source (no incoming): it keeps its own shared value,}\\[2pt]
    \phi_{ij}\big(\sigma_i[\mathrm{sh}\mapsto s_i]\big)
      & \text{$j$ single-source from $i$,}\\[2pt]
    \Gamma_j\big((\sigma_i[\mathrm{sh}\mapsto s_i])_{i \in S(j)}\big)
      & \text{$j$ multi-source.}
  \end{cases}
\]
The morphisms and resolvers are \emph{monotone} when $\Phi$ is monotone:
$\bar{s} \sqsubseteq \bar{s}' \Rightarrow \Phi(\bar{s}) \sqsubseteq \Phi(\bar{s}')$
(componentwise).  Source entries are identities, hence monotone; the
condition constrains only the morphism images and resolvers.  The network is
\emph{valid at bottom} if every component is locally valid with its shared
component set to~$\bot_j$ (for the given locals).
\end{definition}

(Coq: $\Phi$ built from a network, \texttt{fed\_def\_lattice\_shared\_cyclic\_hyps},
which shows that monotone $\Phi$ discharges the hypotheses of the cyclic
normalizer of \texttt{FederationEventsCycles.v}.)

\begin{theorem}[Monotone Convergence Despite Cycles]
\label{thm:monotone-cycles}
Let $\mathcal{N}$ be a morphism network with ordered shared domains
(\Cref{def:lattice-shared}), \emph{possibly cyclic}, whose morphisms and
resolvers are monotone and validity-preserving (M1 / R2), and whose components
satisfy WFC and CC.
\begin{enumerate}[leftmargin=*,nosep,label=(\alph*)]
\item \textbf{Repair normal form.}  $\Phi$ has a least fixed point
  $\bar{s}^{\star}$.  Under the ascending chain condition (ACC; in particular
  for a finite lattice), Kleene iteration from $\bot$ and every fair
  asynchronous (any-order) iteration of the component updates from $\bot$ reach
  $\bar{s}^{\star}$ in finitely many steps, so the repair normal form does not
  depend on the schedule.  Without ACC,
  $\bar{s}^{\star} = \bigsqcup_{k \geq 0} \Phi^k(\bot)$ when $\Phi$ preserves the
  supremum of this chain (for instance, when $\Phi$ is Scott-continuous).
\item \textbf{Validity.}  If moreover the network is valid at bottom, every
  Kleene iterate and every chaotic iterate from $\bot$ is locally valid, and
  under ACC $\bar{s}^{\star}$ is federally valid.  Without ACC the same holds
  when $\Phi$ preserves the supremum of the Kleene chain and local validity is
  closed under suprema of ascending chains.
\item \textbf{Events.}  Let the governed step apply an event and then
  normalize (Phase~1, then reset the shared components to $\bot$ and iterate to
  $\bar{s}^{\star}$).  From a state $\bar{\sigma}^0$, all event sequences that
  differ by swapping federated-independent events reach the same state if and
  only if the \emph{global condition} $\mathrm{GC}(\bar{\sigma}^0)$ holds:
  federated-independent events commute after full re-normalization at every
  state reachable from $\bar{\sigma}^0$.  GC holds at every normalized state
  when, for every target~$j$, with $H_j$ the set of shared values the repair
  writes into~$j$ from valid states (the morphism and resolver images):
  \textbf{(C1$_{\mathrm{cyc}}$)} the local outcome of each event of~$j$ at a
  valid state does not depend on which value of $H_j$ the shared component
  holds; and \textbf{(C2$_{\mathrm{cyc}}$)} for independent events of~$j$, the
  local outcomes of the two orders agree when the shared component is reset to
  the same value of $H_j$ between the events.
\end{enumerate}
\end{theorem}

\begin{proof}
\emph{(a)}~$\Phi$ is a monotone self-map of the complete lattice $\prod_j S_j$ (a
product of complete lattices), so by Knaster--Tarski~\cite{tarski1955} it has a
least fixed point.  Every Kleene iterate lies below every fixed point, by
induction from $\bot$.  Under ACC the ascending chain
$\bot \sqsubseteq \Phi(\bot) \sqsubseteq \cdots$ stabilizes, and a stable value
is a fixed point, hence $\bar{s}^{\star}$.  Order-independence is the
convergence of chaotic iteration~\cite{cousot1977}: below $\bar{s}^{\star}$ a
productive monotone component update strictly increases the state, so under ACC
every fair schedule from $\bot$ terminates, and its result is a fixed point
below $\bar{s}^{\star}$, hence $\bar{s}^{\star}$.  If $\Phi$ preserves the
supremum $s$ of the Kleene chain, then $\Phi(s) = \bigsqcup_k \Phi^{k+1}(\bot)
= s$, and $s$ lies below every fixed point.

\emph{(b)}~M1 and R2 carry local validity from a valid state to its repair, so
by induction from the valid bottom every iterate is locally valid; under ACC
$\bar{s}^{\star}$ is an iterate, and as a fixed point it satisfies every
morphism invariant.  Without ACC, closure of validity under suprema of chains
carries validity to the Kleene supremum.

\emph{(c)}~If GC holds, a swap of two independent events at a reachable state
does not change the result, and trace-equivalent sequences are connected by
such swaps, all at reachable states.  Conversely, if two independent events
fail to commute at a reachable state, the two one-swap sequences through that
state diverge.  For the per-target check: on a valid state the governed step
resets every shared component to $\bot$ and recomputes it from the locals, so
the event's writes to shared components are erased and only its local outcome
matters; C1$_{\mathrm{cyc}}$ and C2$_{\mathrm{cyc}}$ say that this local outcome
does not depend on the order of independent events, since the shared value an
event reads at a normalized state lies in $H_j$. \qed
\end{proof}

(Coq: (a) \texttt{kleene\_lfp}, \texttt{lfp\_unique},
\texttt{kleene\_acc\_lfp}, \texttt{chaotic\_reaches\_lfp},
\texttt{chaotic\_acc\_reaches\_lfp}, \texttt{chaotic\_limit\_unique},
\texttt{kleene\_sup\_lfp}, and for gsm's cyclic normalizer \texttt{cyc\_N\_lfp},
\texttt{cyc\_N\_unique}, \texttt{cyc\_sweep\_order\_independent}; (b)
\texttt{net\_iter\_valid}, \texttt{net\_sweep\_valid}, \texttt{net\_lfp\_valid},
\texttt{net\_Ncyc\_correct}, \texttt{net\_Ncyc\_idem}, \texttt{kleene\_sup\_valid},
non-vacuity \texttt{fed\_valid\_corrected\_instance}; (c) \texttt{gc\_iff},
\texttt{cyc\_events\_converge\_iff}, \texttt{net\_events\_converge\_iff},
\texttt{fed\_thm\_monotone\_cycles\_events\_corrected}, and the per-target check
\texttt{cyc\_check\_gc}, \texttt{cyc\_check\_converges},
\texttt{cyc\_check\_gc\_lfp}, non-vacuity \texttt{cyc\_check\_instance}.)

Each added hypothesis is needed.  Without validity at bottom the least fixed
point need not be federally valid: a 2-cycle of identity morphisms on
$\{\mathit{false} \sqsubset \mathit{true}\}$, valid iff the flag is set, meets
every other hypothesis, yet its least fixed point $(\mathit{false},
\mathit{false})$ is invalid while $(\mathit{true}, \mathit{true})$ is a valid
fixed point (\texttt{fed\_thm\_monotone\_cycles\_lfp\_invalid}).  Without
continuity the Kleene supremum need not be a fixed point: on the complete chain
$0 < 1 < \cdots < \omega < \omega{+}1$ with $\Phi(k) = k{+}1$ and
$\Phi(\omega) = \Phi(\omega{+}1) = \omega{+}1$, the supremum $\omega$ is not
fixed (\texttt{fed\_thm\_monotone\_cycles\_kleene\_formula\_fails}).  Without
GC, events diverge even when every hypothesis of (a) and (b) holds: an event
that copies a shared flag fed around a monotone 2-cycle into a local alarm ends
in different states under the two orders of one allowed swap
(\texttt{cyc\_counterexample}, \texttt{fed\_thm\_monotone\_cycles\_events\_refuted}).
That event fails C1$_{\mathrm{cyc}}$ for every choice of $H_j$
(\texttt{check\_rejects\_latch}); C2$_{\mathrm{cyc}}$ cannot be dropped either
(\texttt{c1\_localcc\_insufficient}), nor replaced by monotonicity of the events
(\texttt{monotone\_c2\_insufficient}).  On a cycle the repair equations can
also have fixed points above the least one (\texttt{bottom\_matters}), which is
why the normalizer resets to $\bot$.

\begin{corollary}[Acyclicity is a concession to non-monotonicity]
\label{cor:acyclicity-monotone}
For the repair normal form, acyclicity is needed only to tame non-monotone
repair.  The cyclic divergence of \Cref{prop:cycle-necessary} uses the antitone
negation $\phi(x) = 1-x$, which is not monotone
(\texttt{negation\_not\_monotone}); under monotonicity
\Cref{thm:monotone-cycles}(a)--(b) applies.  Acyclic structure and monotonicity
are thus two independent routes to a unique federated repair normal form, the
former for arbitrary (possibly non-monotone) repair, the latter for arbitrary
(possibly cyclic) topology.  Neither gives event convergence by itself: that
needs C1 and C2 on acyclic networks (\Cref{thm:fed-convergence,thm:resolved-convergence})
and GC on monotone cycles (\Cref{thm:monotone-cycles}(c)).
\end{corollary}

\begin{remark}[Non-monotone cycles under coordination]
\label{rem:coordinated-cycles}
A non-monotone cycle can still converge if some of its edges are coordinated.
In the invertible fragment (each shared fiber a group acting on itself, each
edge a transport by a group element), pick an authority root and a spanning
tree; drive values along the tree, keep the non-tree edges whose loop is
balanced as checked constraints, and coordinate the unbalanced ones outside the
federation.  The result is the unique state satisfying the tree and the
balanced edges with the root's value, independent of the topological order
(\texttt{coordinated\_sound}, \texttt{coordinated\_unique\_nf}); permutations
of local events converge when the root's own events commute
(\texttt{coordinated\_events\_converge}); and the coordinated set is exactly the
unbalanced edges, no fewer and no more (\texttt{plan\_exact},
\texttt{coordination\_needed}).
\end{remark}

\begin{remark}[Order-independence \emph{is} chaotic iteration]
\Cref{thm:monotone-cycles}(a) exposes the central property of normalization
confluence in another guise: convergence of the repair ``regardless of
application order'' is precisely the order-independent convergence of chaotic
iteration in dataflow analysis and abstract interpretation~\cite{cousot1977}.
The monotone regime inherits that machinery (worklist scheduling) for computing
the federated repair normal form; see \Cref{rem:infinite-lattice} for what
widening does and does not give.
\end{remark}

\begin{remark}[CRDTs are compensation-free monotone federations]
A join-semilattice whose merge is its join is monotone and trivially
validity-preserving, and with every state valid the network is valid at bottom,
so a state-based CRDT~\cite{shapiro2011crdt} fits
\Cref{thm:monotone-cycles}(a)--(b) as the case in which repair never has to
restore an invariant.  (Coq proves CRDT convergence directly,
\texttt{cvrdt\_SEC}, not as an instance of the theorem.)  Normalization
confluence adds compensation on top: monotone convergence \emph{with} business
invariants, which operation commutativity alone does not provide.
\end{remark}

\begin{remark}[Infinite lattices and widening]
\label{rem:infinite-lattice}
\Cref{thm:monotone-cycles} admits a \emph{complete} lattice, which may be
infinite.  Existence and uniqueness of the least fixed point $\bar{s}^{\star}$
hold on any complete lattice (Knaster--Tarski).  Finite iteration reaches it
under ACC.  Without ACC (an unbounded counter or resource in
$\mathbb{N} \cup \{\top\}$, whose ascending chains never stabilize), iteration
may never reach it, even for continuous $\Phi$: on $0 < 1 < \cdots < \omega$
with $\Phi(k) = k{+}1$ and $\Phi(\omega) = \omega$, the least fixed point
$\omega$ is reached by no iterate (\texttt{kleene\_sup\_instance}); and for
non-continuous $\Phi$ not even the limit is the least fixed point
(\Cref{thm:monotone-cycles} and the counterexample after it).  A
\emph{widening} operator (abstract interpretation~\cite{cousot1977}) forces
termination by extrapolating to a post-fixed point $w$, $\Phi(w) \sqsubseteq w$,
which bounds every iterate and the least fixed point
(\texttt{widening\_sound}).  In general $w$ is not a fixed point, so it does not
satisfy the morphism invariants and is not the federated normal form
(\texttt{widening\_not\_normal\_form}).  Widening gives a sound bound, not
convergence.
\end{remark}

\subsection{Compositionality: Sub-Federations as Effective Registries}
\label{sec:compositionality}

Product composition (\Cref{thm:product}) collapses \emph{independent} registries
into one.  The same holds for a \emph{connected} sub-federation: an internally
convergent sub-network can be replaced by a single \emph{effective registry},
so federations verify modularly and hierarchically.  Registries and morphisms
form a category in which the federated registry is the limit of the constraint
diagram; this subsection is the statement that such limits compose.

\begin{definition}[Sub-Federation, Boundary, Convexity]
\label{def:subfed}
Given $\mathcal{N} = (\{R_i\}_{i \in \mathcal{I}}, \{\phi_{ij}\})$ and
$\mathcal{J} \subseteq \mathcal{I}$, the \emph{sub-federation} on $\mathcal{J}$
has components $\{R_i\}_{i \in \mathcal{J}}$ and \emph{internal} morphisms
$E_{\mathcal{J}} = \{(i,j) \in E : i, j \in \mathcal{J}\}$.  Its boundary
consists of \emph{import} edges (from $\mathcal{I}\setminus\mathcal{J}$ into
$\mathcal{J}$) and \emph{export} edges (from $\mathcal{J}$ out).  $\mathcal{J}$
is \emph{convex} if no directed path between two members of $\mathcal{J}$ passes
through a vertex outside $\mathcal{J}$.
\end{definition}

\begin{definition}[Effective Registry]
\label{def:effreg}
The \emph{effective registry} $R_{\mathcal{J}}$ is the federated registry
(\Cref{def:fed-registry}) of the sub-federation: state space
$\prod_{i \in \mathcal{J}} \Sigma_{R_i}$, invariants the internal local and
morphism invariants, and compensation the sub-federated normalizer
$\rho_{\Fed}^{\mathcal{J}}$.  Its shared component is the union of the import
ports (the shared components of members targeted by import edges); export edges
read its normalized state.
\end{definition}

\begin{theorem}[Compositional Collapse]
\label{thm:collapse}
Let $\mathcal{J}$ be convex with an internally convergent sub-federation: its
morphisms satisfy M1/R2, its resolvers R1, its components WFC and CC, and
either (i) it is acyclic and satisfies C1 and C2 (\Cref{def:c1c2}) on its
internal edges, or (ii) it is cyclic with monotone repair under ACC, valid at
bottom, and satisfies C1$_{\mathrm{cyc}}$ and C2$_{\mathrm{cyc}}$
(\Cref{thm:monotone-cycles}).  Then:
\begin{enumerate}[leftmargin=*,nosep,label=(\alph*)]
\item $R_{\mathcal{J}}$ satisfies WFC, and its governed step
  $\bar{\sigma} \mapsto \rho_{\Fed}^{\mathcal{J}}(\apply(e, \bar{\sigma}))$
  maps valid states to valid states and satisfies CC1 at every valid state of
  $R_{\mathcal{J}}$ (the component condition the guarded theorems use).
\item Each boundary morphism, reinterpreted as a morphism to or from
  $R_{\mathcal{J}}$, satisfies M1 iff it does in $\mathcal{N}$.
\item Let $\mathcal{N}'$ replace the sub-federation by $R_{\mathcal{J}}$.  Then
  $\mathcal{N}'$ is a well-formed network, acyclic if $\mathcal{N}$ is;
  $\rho_\Fed$ on $\mathcal{N}$ and on $\mathcal{N}'$ agree under
  $\Sigma_{R_{\mathcal{J}}} = \prod_{i \in \mathcal{J}} \Sigma_{R_i}$, and so do
  their governed steps.  Hence the two networks have the same runs from every
  federally valid state: $\mathcal{N}$ meets the exact condition of
  \Cref{thm:fed-convergence}(b) iff $\mathcal{N}'$ does, and their normal forms
  agree.
\end{enumerate}
\end{theorem}

\begin{proof}
\emph{(a)}~WFC: $\rho_{\Fed}^{\mathcal{J}}$ reaches a federally valid state of
the sub-federation in one application (\Cref{lem:fed-termination},
\Cref{lem:resolved-termination}; in case~(ii) \Cref{thm:monotone-cycles}(b)), so
the measure $\mathbf{1}[\text{invalid}]$ is a well-founded compensation measure
for $R_{\mathcal{J}}$.  CC1 at valid states: in case~(i) this is
\Cref{thm:fed-cc} applied inside $\mathcal{J}$, with the import ports as
additional sources whose values are held fixed; in case~(ii) it is GC at
normalized states, \Cref{thm:monotone-cycles}(c).

\emph{(b)}~An import edge $(i,j)$ with $j \in \mathcal{J}$ overwrites the shared
component of $R_{\mathcal{J}}$ that coincides with $\pi_j^{\mathrm{sh}}$; its M1
obligation (overwriting with a valid source image preserves target validity) is
unchanged, since sub-federated re-normalization preserves federal validity
by~(a).  An export edge reads $R_{\mathcal{J}}$'s normalized state, a
source-determined projection.

\emph{(c)}~Convexity ensures contracting $\mathcal{J}$ to a single vertex creates
no exterior path between two members, so $\mathcal{N}'$ is acyclic whenever
$\mathcal{N}$ is.  Because $\mathcal{J}$ is convex, a topological order of
$\mathcal{N}$ visits it as a contiguous block, so $\rho_{\Fed}$ on $\mathcal{N}$
factors as ``normalize the sub-federation, then propagate across the boundary'',
which is exactly $\rho_{\Fed}$ on $\mathcal{N}'$ with $R_{\mathcal{J}}$'s
compensation inlined (the fold-append law of the staged normalizer).  An event
of $\mathcal{J}$ changes one member and is followed by the same repair on both
sides, so the governed steps agree, and so do all runs. \qed
\end{proof}

This theorem is proved on paper.  Its operator-level core, the fold-append law
for the staged normalizer, is mechanized (\texttt{rhoFold\_compositional},
\texttt{rhoF\_from\_app}); the effective registry, the contraction of a convex
block, and parts (a)--(c) are not yet.  Version~1 stated (a) under component CC
alone, which relied on the refuted federated CC theorem; the hypotheses above
add C1 and C2 (or the cyclic checks).

\begin{corollary}[Modular and Hierarchical Verification]
\label{cor:modular}
A federation verifies compositionally: verify each convex sub-federation once
(its components, M1/R1/R2, and C1 and C2 on its internal edges), replace it by
its effective registry, and verify the network of effective registries (M1/R2
and C1 and C2 on the remaining edges).  The construction nests, giving
hierarchical verification, and a verified sub-federation is reusable as a
black-box component with a fixed port interface.  \Cref{thm:product} is the
special case $E_{\mathcal{J}} = \varnothing$ with empty boundary (independent
registries).  This corollary inherits the paper-level status of
\Cref{thm:collapse}.
\end{corollary}

\begin{remark}[Convexity and the monotone regime]
Convexity is used to keep the collapsed network acyclic.  In the monotone
regime a boundary-induced cycle is harmless provided $\rho_{\Fed}^{\mathcal{J}}$
is monotone, and then \Cref{thm:monotone-cycles} applies to $\mathcal{N}'$
(with validity at bottom for~(b) and GC for events).  Monotone repair inside
$\mathcal{J}$ is \emph{not} enough for this: $\rho_{\Fed}^{\mathcal{J}}$ begins
with the components' own normalizers (Phase~1), which need not be monotone.  A
one-registry sub-federation with no internal edges has constant, hence
monotone, repair, yet its normalizer sends $a_0 \sqsubseteq a_1$ to
$a_2 \not\sqsubseteq a_1$ (\texttt{fed\_rem\_convexity\_refuted}).  If Phase~1
is monotone and $\Phi$ is monotone in the locals as well as in the shared
values, then under ACC $\rho_{\Fed}^{\mathcal{J}}$ is monotone
(\texttt{fed\_rem\_convexity\_corrected}, \texttt{lfp\_mono\_param}; non-vacuity
\texttt{convexity\_corrected\_instance}).
\end{remark}

\subsection{Example: Manufacturer--Supplier Federation}

\textbf{Manufacturer $R_M$}: $\Sigma_M = \{\texttt{draft}, \texttt{active},
\texttt{suspended}\}$, all locally valid.
$e_{\mathrm{pub}}$: $\texttt{draft} \to \texttt{active}$.

\textbf{Supplier $R_S$}: $\Sigma_S = \{\texttt{idle}, \texttt{listed},
\texttt{stale}, \texttt{err}\}$; $V_S(\texttt{err}) = \bot$, rest valid;
$\rho_S(\texttt{err}) = \texttt{idle}$.
$e_{\mathrm{exp}}$: sets state to $\texttt{stale}$.

\textbf{Total morphism $\phi$}:
$\texttt{draft} \mapsto \texttt{idle}$,
$\texttt{active} \mapsto \texttt{listed}$,
$\texttt{suspended} \mapsto \texttt{stale}$.
M1 holds (all images are valid).

\textbf{Convergence trace} from $(\texttt{draft}, \texttt{idle})$ with
independent events $e_{\mathrm{pub}}$ (targets $R_M$) and $e_{\mathrm{exp}}$
(targets $R_S$):

\emph{Order 1} ($e_{\mathrm{pub}}$ first):
$(\texttt{active}, \texttt{idle}) \xrightarrow{\phi}
(\texttt{active}, \texttt{listed}) \xrightarrow{e_{\mathrm{exp}}}
(\texttt{active}, \texttt{stale}) \xrightarrow{\phi}
(\texttt{active}, \texttt{listed})$.

\emph{Order 2} ($e_{\mathrm{exp}}$ first):
$(\texttt{draft}, \texttt{stale}) \xrightarrow{\phi}
(\texttt{draft}, \texttt{idle}) \xrightarrow{e_{\mathrm{pub}}}
(\texttt{active}, \texttt{idle}) \xrightarrow{\phi}
(\texttt{active}, \texttt{listed})$.

Both converge to $(\texttt{active}, \texttt{listed})$.  The supplier event's
effect is overwritten by morphism repair: the manufacturer is authoritative.
This is an instance of \Cref{thm:fed-convergence}, not a coincidence of the two
traces: M1, the supplier's WFC, C1 and C2 all hold (the supplier's event writes
only state that the final repair overwrites, so both sides of C1 agree), and
every reordering of any event list converges from every federally valid state
(Coq: \texttt{ms\_instance}, \texttt{ms\_paper\_traces}, \texttt{ms\_converges}).

\begin{remark}[Machine-checked]
\label{rem:fed-mechanized}
The federated results of this section are mechanized in Coq/Rocq, in the
\texttt{coq/} directory of the artifact, against the paper's own hypotheses:
\texttt{FederationGRS.v} states (H) and (H$_R$) as a record (shared/local
decomposition, acyclic network, M1, R1, R2, component WFC with $\rho^*$ as the
iterate, component CC1 and CC2) and proves every result above under its label,
for the guarded system $\mathcal{G}_\Fed$ and for the unguarded one.  The
machine model of a federated processor and the sufficiency of C1 and C2 are in
\texttt{FederationEvents.v}; the exact converses in
\texttt{FederationEventsConverse.v}; monotone cycles in
\texttt{FederationEventsCycles.v}, \texttt{FederationEventsCyclesCheck.v},
\texttt{MonotoneFederation.v}, \texttt{Chaotic.v} and \texttt{ChaoticACC.v};
coordinated cycles in \texttt{CoordinatedCycles.v}; and the concrete examples
(\Cref{prop:cycle-necessary,prop:m1-necessary}, the resolution conditions, and
the manufacturer and supplier) in \texttt{PaperInstances.v}.  Each refuted
statement of version~1 has a gated counterexample against its exact hypotheses
(Appendix~\ref{app:errata}).  The theorems named in this paper are in the
\texttt{Print Assumptions} gate of \texttt{coq/verify.sh} (GATECOUNT theorems
at the time of writing), and each is ``Closed under the global context'': no
axioms, no admitted lemmas.  Not mechanized: the Compositional Collapse
(\Cref{thm:collapse}, \Cref{cor:modular}) beyond its fold-append core, and the
decidability remark for R1, R2, C1 and C2.
\end{remark}

%========================================
\section{Discussion}
\label{sec:discussion}

\subsection{Three Regimes of Coordination-Free Convergence}
\label{sec:design-space}

Our result, combined with CRDTs~\cite{shapiro2011crdt} and invariant
confluence~\cite{bailis2014coordination}, identifies three distinct structural
regimes under which distributed processors converge without coordination:

\begin{enumerate}[leftmargin=*,nosep]
\item \textbf{Operation commutativity} (CRDTs): Operations commute
  algebraically.  Convergence follows from the algebra of operations.  No
  invariant enforcement or compensation is needed.
\item \textbf{Invariant confluence} (I-confluence): Operations preserve
  application invariants under all orderings.  Convergence follows because
  validity is never violated.  No compensation is needed.
\item \textbf{Normalization confluence}~\cite{blackwell2026nc}: Operations may be
  non-commutative and may violate invariants.  Compensation normalizes states
  to validity.  WFC and CC are sufficient for semantic convergence
  (\Cref{thm:convergence}); UBC is orthogonal to convergence and governs
  scalability (\Cref{thm:complexity}).  No coordination is needed.
\end{enumerate}

The three regimes form a containment hierarchy.  Operation commutativity
implies CC (with $\rho_R^* = \mathrm{id}$).  Invariant confluence implies that
compensation never fires.  Normalization confluence subsumes both as special
cases and additionally covers systems where operations are non-commutative and
invariant-violating---the common case in business applications.

The regimes partition the coordination-free design space by \emph{where}
convergence is guaranteed: in the operations (regime~1), in the invariants
(regime~2), or in the normalization rewrite system (regime~3).

\begin{remark}[Convergence as Rewrite Property]
\label{rem:rewrite-property}
The conventional literature frames convergence as a property of operations:
operations must commute, or operations must preserve invariants.
Normalization confluence reframes convergence as a property of the
\emph{rewrite system} that interleaves application and compensation.  The
operations themselves need no special algebraic structure.  This shifts the
design burden from operation algebra to compensation design, and places the
convergence mechanism at the semantic normalization layer---above the
data-structure layer (CRDTs) and the transaction layer (I-confluence).
Moreover, this rewrite-system perspective extends across organizational
boundaries: in the federated results (\Cref{sec:federation}) multi-registry
systems inherit termination and a unique repair normal form from their
components, with morphism repair as additional rewrite rules in the same
framework; event convergence additionally needs two per-edge checks, C1 and C2.
\end{remark}

\subsection{Two Regimes for Federated Repair}
\label{sec:regimes}

The federated results partition constraint-network design into two regimes with
\emph{independent} routes to a unique repair normal form
(\Cref{cor:acyclicity-monotone}).  When repair is \emph{monotone} on ordered
shared domains, the repair normal form exists on \emph{any} topology, including
cyclic meshes, by \Cref{thm:monotone-cycles} (federally valid when the network
is valid at bottom).  When repair may be \emph{non-monotone} (toggles,
cancellations, order-reversing revocations), the topology must be
\emph{acyclic}, with resolution operators (\Cref{sec:multisource}) at
multi-source targets, and the normal form is the one-shot constructive one
(\Cref{cor:resolved-nf}), unless the offending cycles are coordinated
(\Cref{rem:coordinated-cycles}).  In both regimes repair must preserve validity
(M1/R2), and in both, convergence of \emph{event interleavings} needs a further
check: C1 and C2 per edge on acyclic networks, GC on monotone cycles (implied by
per-target C1 and C2 over the repair images).  \Cref{tab:regimes} summarizes
the trade.

\begin{table}[ht]
\centering
\small
\begin{tabular}{@{}lp{3.0cm}p{2.6cm}p{3.0cm}p{2.6cm}@{}}
\toprule
\textbf{Regime} & \textbf{Repair requirement} & \textbf{Topologies} & \textbf{Repair normal form} & \textbf{Event convergence} \\
\midrule
Monotone &
Monotone on a lattice order (join \emph{or} meet); validity-preserving; valid at bottom &
Any: cyclic meshes, DAGs, trees &
Least fixed point $\mathrm{lfp}(\bot)$ under any \emph{fair} asynchronous order (Knaster--Tarski, chaotic iteration, ACC) &
Iff GC; implied by per-target C1, C2 over images \\[4pt]
Non-monotone &
May reverse the order (toggles, negation, cancellations); validity-preserving &
Acyclic only; multi-source targets need resolution operators &
Deterministic one-shot constructive normal form &
Under C1, C2; exact at reachable witnesses \\
\bottomrule
\end{tabular}
\caption{Two regimes for federated repair.  Monotonicity and acyclicity are
independent conditions for a unique repair normal form: monotone repair on any
topology (\Cref{thm:monotone-cycles}); non-monotone repair requires acyclicity
(\Cref{prop:cycle-necessary}).  Validity preservation (M1/R2) is required in
both, and event convergence needs the checks in the last column
(\Cref{thm:fed-convergence,thm:resolved-convergence,thm:monotone-cycles}).}
\label{tab:regimes}
\end{table}

\subsection{Verifying Compensation Commutativity}

CC is easy to state but potentially hard to verify: it quantifies over all
states and all independent event pairs.  \Cref{sec:calculus} develops a
verification calculus that reduces CC to local, per-event-pair checks via
invariant footprints and decomposable repair.  Here we note three common
patterns that practitioners can use as a first check before invoking the full
calculus:

\begin{enumerate}[leftmargin=*,nosep]
\item \textbf{Compensation to a canonical state.}\;  If there exists a
  default valid state $\sigma_\bot$ such that $\rho_R^*(\sigma) = \sigma_\bot$
  for all invalid~$\sigma$, both CC1 and CC2 hold trivially.  This is the
  ``reset to safe default'' pattern.
\item \textbf{Lattice-structured compensation.}\;  If the valid states form a
  join-semilattice and compensation computes the join with the nearest valid
  state, both conditions hold by the commutativity of the join operation.
\item \textbf{Decomposable repair.}\;  If compensation decomposes into
  commuting per-invariant repairs, CC can be verified via the footprint
  calculus of \Cref{sec:calculus}.  The order fulfillment example
  (\Cref{sec:example}) uses this pattern.
\end{enumerate}

\subsection{Applications}

The compensate-and-converge principle applies wherever systems process streams
of changes under semantic constraints with bounded compensation.  We highlight
three domains.

\paragraph{Infrastructure Configuration Drift.}
Infrastructure-as-code tools apply streams of configuration changes to
infrastructure state.  Valid configurations are defined by policy.  The
semantic state space is finite ($k$ boolean policy predicates yield at most
$2^k$ equivalence classes).  Remediation operators that reset violated policies
to compliant defaults satisfy CC2 (compensation to a canonical state per
policy).  Our theorem guarantees convergence to equivalent policy-compliant
states regardless of the order in which changes are applied.

\paragraph{Database Schema Migration.}
Schema migrations are event streams applied to schema state.  The set of valid
schema configurations is finite.  Migration scripts that repair constraint
violations constitute compensation operators.  When such scripts resolve each
constraint independently (lattice-structured compensation), CC holds and our
theorem guarantees order-independent convergence.

\paragraph{Regulatory Compliance.}
Organizations processing streams of regulatory changes must maintain
compliance.  Compliance states are finite (each regulation is either satisfied
or not).  Remediation procedures that restore compliance by addressing each
violated regulation independently satisfy CC.  Our theorem guarantees
convergence to equivalent compliance states.

\paragraph{Cross-Organizational Supply Chains.}
The federated convergence result (\Cref{sec:federation}) applies to
multi-registry settings.  A manufacturer's product specification registry
connected by morphisms to supplier capability registries forms a tree-shaped
network.  Each organization maintains its own compensation operators; the
morphisms encode specification conformance.  When each supplier's events pass
C1 and C2 against the manufacturer's images (a per-edge check),
\Cref{thm:fed-convergence} guarantees that all participants converge to
consistent states across organizational boundaries.  When a supplier answers to \emph{several}
authorities at once---say a manufacturer's specifications \emph{and} a
regulator's rules---the network is no longer a tree; a resolution operator
(\Cref{sec:multisource}) combines the constraints (for instance, most-
restrictive-wins), and \Cref{thm:resolved-convergence} extends the same
guarantee to this acyclic multi-source setting.

\subsection{Extensions}

\paragraph{At-least-once delivery.}
The convergence theorems assume each event is delivered once.  Real transports
deliver at least once, so an event may be redelivered, possibly much later.  A
duplicate is absorbed when the redelivered event's governed step is idempotent
and commutes with every event delivered between its copies (Coq:
\texttt{alo\_absorbed}); if all events commute and every duplicated event is
idempotent, any two at-least-once deliveries of the same events agree
(\texttt{alo\_commuting\_converges}); under causal delivery only concurrent
events need to commute, provided redelivery is itself causally consistent
(\texttt{causal\_alo\_exactly\_once}, \texttt{causal\_alo\_converges}).
Both qualifiers are needed: a duplicate of an event whose governed step is not
idempotent diverges (\texttt{non\_idempotent\_diverges}), and a late
redelivery of an idempotent event after its causal successor can diverge
(\texttt{late\_duplicate\_diverges}).

\paragraph{Privacy.}
One extension remains future work: \emph{privacy-preserving validation} where
organizations prove registry compliance via zero-knowledge proofs without
revealing state.  The federated convergence result (\Cref{sec:federation})
handles the multi-registry coordination problem; privacy-preserving validation
would additionally protect organizational sovereignty over state visibility.

%========================================
\section{Verification Calculus for CC}
\label{sec:calculus}

Compensation commutativity is a global condition: it quantifies over all states
and all independent event pairs.  This section, based on~\cite{blackwell2026nc},
develops sufficient conditions that reduce CC to local, per-event-pair checks.
The key tools are an \emph{invariant footprint} mapping and a
\emph{decomposable repair} condition.

\subsection{Notation}

For readability in algebraic manipulations, we write
$N(\sigma) \defeq \rho_R^*(\sigma)$ for iterated compensation to validity
(\Cref{def:rhostar}).  The normalizer satisfies: (i)~$V_R(N(\sigma)) = \top$
for all~$\sigma$; (ii)~$N(\sigma) = \sigma$ if $V_R(\sigma) = \top$;
(iii)~$N(N(\sigma)) = N(\sigma)$ (idempotence).  We write
$A_e(\sigma) \defeq \mathrm{apply}(e, \sigma)$ when convenient.

In this notation, the CC obligations become:
\begin{itemize}[leftmargin=*,nosep]
\item[\textbf{(CC1)}] $N(A_{e_2}(N(A_{e_1}(\sigma)))) =
  N(A_{e_1}(N(A_{e_2}(\sigma))))$ for independent $e_1, e_2$.
\item[\textbf{(CC2)}] $N(A_e(\sigma)) = N(A_e(\rho_R(\sigma)))$ for
  invalid~$\sigma$.
\end{itemize}

\subsection{Invariant Vectors and Footprints}

\begin{definition}[Invariant Vector]
\label{def:inv-vector}
Define $\mathbf{v}: \Sigma_R \to \{0,1\}^k$ by
$\mathbf{v}(\sigma)_i = 1 \iff \psi_i(\sigma) = \top$.
\end{definition}

\begin{definition}[Footprint]
\label{def:footprint}
A \emph{footprint} is a function $F: \mathcal{E} \to 2^{\{1,\ldots,k\}}$ such
that for all $\sigma$ and all $i \notin F(e)$,
\[
\mathbf{v}(A_e(\sigma))_i = \mathbf{v}(\sigma)_i.
\]
Intuitively, $F(e)$ is the set of invariants whose truth value can change
under~$e$.  Over-approximations are sound: larger footprints make the calculus
more conservative but never incorrect.
\end{definition}

\subsection{Decomposable Repair}

We capture the common case where repair addresses each violated invariant
independently.

\begin{definition}[Per-Invariant Repair]
\label{def:perinv}
A set of \emph{per-invariant repair operators} $r_1, \ldots, r_k :
\Sigma_R \to \Sigma_R$ satisfies:
\begin{enumerate}[leftmargin=*,nosep,label=(R\arabic*)]
\item \textbf{Targeted fix:}\; if $\psi_i(\sigma) = \bot$, then
  $\psi_i(r_i(\sigma)) = \top$.
\item \textbf{No-op on satisfied:}\; if $\psi_i(\sigma) = \top$, then
  $r_i(\sigma) = \sigma$.
\item \textbf{Mutual commutativity:}\; $r_i(r_j(\sigma)) = r_j(r_i(\sigma))$
  for all $i, j$.
\end{enumerate}
\end{definition}

\begin{definition}[Decomposable Normalizer]
\label{def:decomp}
The normalizer $N$ is \emph{decomposable} if there exist per-invariant repairs
$r_1, \ldots, r_k$ as above such that
\[
N(\sigma) = r_1(r_2(\cdots r_k(\sigma) \cdots))
\]
for all~$\sigma$.  (The order is irrelevant by mutual commutativity.)
\end{definition}

\begin{definition}[Repair Locality]
\label{def:repair-locality}
An event $e$ is \emph{repair-local} with respect to footprint $F(e)$ if for
all $i \notin F(e)$,
\[
r_i(A_e(\sigma)) = A_e(r_i(\sigma)).
\]
That is, $e$ commutes with repairs of invariants it does not affect.
\end{definition}

\subsection{Strong Absorption Implies CC2}

\begin{theorem}[Strong Absorption]
\label{thm:strong-absorption}
Assume WFC and let $N = \rho_R^*$.  If for all events $e$ and states~$\sigma$,
\[
N(A_e(\sigma)) = N(A_e(N(\sigma))),
\]
then CC2 holds.
\end{theorem}

\begin{proof}
Let $V_R(\sigma) = \bot$.  Because $N$ iterates $\rho_R$ to the unique valid
fixpoint reachable from~$\sigma$, and $\rho_R(\sigma)$ lies on the same
compensation chain, we have $N(\sigma) = N(\rho_R(\sigma))$.

Applying the assumed equation at $\sigma$ and at $\rho_R(\sigma)$:
\[
N(A_e(\sigma)) = N(A_e(N(\sigma))) = N(A_e(N(\rho_R(\sigma)))) = N(A_e(\rho_R(\sigma))),
\]
where the first and third equalities use the assumption and the middle equality
uses $N(\sigma) = N(\rho_R(\sigma))$.  This is CC2.
\end{proof}

\begin{remark}
Strong absorption states: ``normalizing before applying an event does not
change the final normalized outcome.''  This holds whenever compensation
repairs violations in a way determined by the violation pattern rather than by
the path to the violation.  The order fulfillment example
(\Cref{sec:example}) satisfies this: the held state captures the approval
intent regardless of whether the order was directly approved with low balance
or arrived at the invalid state through a different path.
\end{remark}

\subsection{Footprint Disjointness Implies CC1}

We prove two preparatory lemmas before the main theorem.

\begin{lemma}[Repair-Event Commutation]
\label{lem:repair-commute}
If event $e$ is repair-local with respect to footprint $F(e)$
(\Cref{def:repair-locality}), then for any index set $J$ with
$J \cap F(e) = \varnothing$ and any state~$\sigma$:
\[
R_J(A_e(\sigma)) = A_e(R_J(\sigma)),
\]
where $R_J = \prod_{i \in J} r_i$.
\end{lemma}

\begin{proof}
By repair locality, $r_i \circ A_e = A_e \circ r_i$ for each $i \in J$
(since $i \notin F(e)$).  The result follows by induction on $|J|$: if the
claim holds for $R_{J'}$ with $|J'| < |J|$, then for any $j \in J$,
\[
R_J(A_e(\sigma)) = r_j(R_{J \setminus \{j\}}(A_e(\sigma)))
= r_j(A_e(R_{J \setminus \{j\}}(\sigma)))
= A_e(r_j(R_{J \setminus \{j\}}(\sigma)))
= A_e(R_J(\sigma)),
\]
using the inductive hypothesis and then $r_j \circ A_e = A_e \circ r_j$.
\end{proof}

\begin{lemma}[Repair Block Idempotence]
\label{lem:repair-idempotent}
For any index set $J \subseteq \{1,\ldots,k\}$, the repair block $R_J$ is
idempotent on states that are valid on all invariants in~$J$: if
$\psi_i(\sigma) = \top$ for all $i \in J$, then $R_J(\sigma) = \sigma$.
More generally, $R_J(R_J(\sigma)) = R_J(\sigma)$ for all~$\sigma$.
\end{lemma}

\begin{proof}
If $\psi_i(\sigma) = \top$ for all $i \in J$, then each $r_i(\sigma) = \sigma$
by~(R2), so $R_J(\sigma) = \sigma$.  For general idempotence: $R_J(\sigma)$
satisfies $\psi_i(R_J(\sigma)) = \top$ for all $i \in J$ (each $r_i$ fixes its
invariant by~(R1) and does not break others by mutual commutativity and~(R2)),
so $R_J(R_J(\sigma)) = R_J(\sigma)$ by the first claim.
\end{proof}

\begin{theorem}[Footprint Commutativity]
\label{thm:footprint-cc1}
Assume WFC and that $N$ is decomposable via commuting per-invariant repairs
$r_1, \ldots, r_k$.  Let $e_1, e_2$ be independent events satisfying:
\begin{enumerate}[leftmargin=*,nosep,label=(H\arabic*)]
\item Disjoint footprints: $F(e_1) \cap F(e_2) = \varnothing$.
\item Repair locality: each $e_j$ is repair-local with respect to $F(e_j)$.
\end{enumerate}
Then CC1 holds for $(e_1, e_2)$.
\end{theorem}

\begin{proof}
Partition $\{1, \ldots, k\}$ into three disjoint index sets:
$S_1 = F(e_1)$, $S_2 = F(e_2)$,
$S_0 = \{1,\ldots,k\} \setminus (S_1 \cup S_2)$.
By decomposability and mutual commutativity, $N = R_{S_0} \circ R_{S_1}
\circ R_{S_2}$ in any order.  We compute both sides of CC1 as chains of
equalities.

\emph{Left side}: $N(A_{e_2}(N(A_{e_1}(\sigma))))$.
\begin{align}
&= N(A_{e_2}(R_{S_0}(R_{S_1}(R_{S_2}(A_{e_1}(\sigma))))))
  \tag{expand inner $N$} \\
&= N(R_{S_0}(R_{S_1}(A_{e_2}(R_{S_2}(A_{e_1}(\sigma))))))
  \tag{\Cref{lem:repair-commute}: $S_0 \cap F(e_2) = S_1 \cap F(e_2) = \varnothing$} \\
&= R_{S_0}(R_{S_1}(R_{S_2}(R_{S_0}(R_{S_1}(A_{e_2}(R_{S_2}(A_{e_1}(\sigma))))))))
  \tag{expand outer $N$} \\
&= R_{S_0}(R_{S_1}(R_{S_2}(A_{e_2}(R_{S_2}(A_{e_1}(\sigma))))))
  \tag{\Cref{lem:repair-idempotent}: $R_{S_0}, R_{S_1}$ absorbed}
\end{align}

\emph{Right side}: $N(A_{e_1}(N(A_{e_2}(\sigma))))$.
\begin{align}
&= N(A_{e_1}(R_{S_0}(R_{S_2}(R_{S_1}(A_{e_2}(\sigma))))))
  \tag{expand inner $N$} \\
&= N(R_{S_0}(R_{S_2}(A_{e_1}(R_{S_1}(A_{e_2}(\sigma))))))
  \tag{\Cref{lem:repair-commute}: $S_0 \cap F(e_1) = S_2 \cap F(e_1) = \varnothing$} \\
&= R_{S_0}(R_{S_1}(R_{S_2}(R_{S_0}(R_{S_2}(A_{e_1}(R_{S_1}(A_{e_2}(\sigma))))))))
  \tag{expand outer $N$} \\
&= R_{S_0}(R_{S_1}(R_{S_2}(A_{e_1}(R_{S_1}(A_{e_2}(\sigma))))))
  \tag{\Cref{lem:repair-idempotent}: $R_{S_0}, R_{S_2}$ absorbed}
\end{align}

\emph{Equating.}  The left side is
$R_{S_0}(R_{S_1}(R_{S_2}(A_{e_2}(R_{S_2}(A_{e_1}(\sigma))))))$ and the right
side is $R_{S_0}(R_{S_1}(R_{S_2}(A_{e_1}(R_{S_1}(A_{e_2}(\sigma))))))$.
It remains to show their inner expressions are equal.

Apply \Cref{lem:repair-commute} to move $R_{S_2}$ past $A_{e_1}$ on the left
(since $S_2 \cap F(e_1) = \varnothing$):
\[
A_{e_2}(R_{S_2}(A_{e_1}(\sigma))) = A_{e_2}(A_{e_1}(R_{S_2}(\sigma))).
\]
Apply \Cref{lem:repair-commute} to move $R_{S_1}$ past $A_{e_2}$ on the right
(since $S_1 \cap F(e_2) = \varnothing$):
\[
A_{e_1}(R_{S_1}(A_{e_2}(\sigma))) = A_{e_1}(A_{e_2}(R_{S_1}(\sigma))).
\]
Both inner expressions now have the form
$A_{e_j}(A_{e_\ell}(R_?(\sigma)))$ with all repairs pushed past the events.
Applying $R_{S_2}$ to the left expression and $R_{S_1}$ to the right, then
using mutual commutativity of $R_{S_1}$ and $R_{S_2}$, both sides reduce to
\[
R_{S_0}(R_{S_1}(R_{S_2}(A_{e_2}(A_{e_1}(R_{S_2}(R_{S_1}(\sigma))))))).
\]
Since $R_{S_1}$ and $R_{S_2}$ commute and each is idempotent after application
(\Cref{lem:repair-idempotent}), the left and right sides are equal.
\end{proof}

\begin{remark}[Practical import]
A practitioner does not prove CC1 by global state enumeration.  Instead:
(i)~assign each event type an invariant footprint, (ii)~verify that each event
commutes with repairs outside its footprint (\Cref{lem:repair-commute}),
(iii)~check that independent event pairs have disjoint footprints.  All three
are local obligations.  The order fulfillment example (\Cref{sec:example})
demonstrates this workflow: $e_{\mathrm{credit}}$ and $e_{\mathrm{approve}}$
affect disjoint aspects of the state, and the held-state repair commutes with
credit application.
\end{remark}

\subsection{Product Composition}

Real systems comprise multiple semantic modules.  The following theorem shows
CC verifications compose.

\begin{definition}[Product Registry]
\label{def:product}
Let $R_1 = (\Sigma_1, I_1, V_1, \rho_1)$ and $R_2 = (\Sigma_2, I_2, V_2,
\rho_2)$ with disjoint invariant sets.  Define $R_1 \otimes R_2$ by:
\[
\Sigma = \Sigma_1 \times \Sigma_2, \quad
V(\sigma_1, \sigma_2) = V_1(\sigma_1) \wedge V_2(\sigma_2), \quad
\rho(\sigma_1, \sigma_2) = (\rho_1(\sigma_1), \rho_2(\sigma_2)).
\]
Events are pairs $e = (e^{(1)}, e^{(2)})$ with application componentwise.
\end{definition}

\begin{theorem}[Product Lifting]
\label{thm:product}
If each $R_j$ satisfies WFC and CC, then $R_1 \otimes R_2$ satisfies WFC and
CC.  If each normalizer is decomposable, the product normalizer is
decomposable with footprints composed by union across components.
\end{theorem}

\begin{proof}
\emph{WFC}: Use measure $\Phi(\sigma_1, \sigma_2) = \Phi_1(\sigma_1) +
\Phi_2(\sigma_2)$.  Each compensation step strictly decreases at least one
component when invalid.

\emph{CC2}: The product normalizer satisfies $N(\sigma_1, \sigma_2) =
(N_1(\sigma_1), N_2(\sigma_2))$ and application is componentwise, so
CC2 holds iff it holds in each component.

\emph{CC1}: Both sides of the CC1 equation decompose into paired equalities
in $R_1$ and $R_2$ by componentwise application and normalization.

\emph{Decomposability}: Per-invariant repairs in each component act on
disjoint state components and thus commute across components trivially.
\end{proof}

\begin{remark}
Product composition enables modular system design: verify CC for each semantic
module independently, then compose without re-verification.  This handles the
\emph{independent} case (disjoint invariants, no cross-module constraints).
The federated convergence result (\Cref{sec:federation}) generalizes product
composition to the \emph{dependent} case: registries connected by morphisms
encoding cross-organizational semantic constraints.
\end{remark}

\subsection{Partial Decomposability}

Full decomposability is a strong sufficient condition.  When it does not hold
globally, the calculus still provides value:

\begin{enumerate}[leftmargin=*,nosep]
\item CC may hold for event pairs with disjoint footprints even if the full
  normalizer is not decomposable.  The footprint theorem
  (\Cref{thm:footprint-cc1}) applies to any event pair satisfying its
  hypotheses, regardless of whether \emph{all} pairs do.
\item When footprints overlap for some event pair, the calculus identifies
  precisely which pairs require further analysis.  A practitioner can then
  either redesign compensation to be canonical for the overlapping invariants
  or add coordination only for the interacting event types.
\item For event types that share footprints but whose normalized outcomes
  agree by domain-specific reasoning (e.g., both repairs clamp to the same
  bound), CC1 can be verified by case analysis restricted to the overlapping
  states.  The footprint calculus reduces the case space.
\end{enumerate}

%========================================
\section{Conclusion}
\label{sec:conclusion}

We have identified \emph{normalization confluence} as a third structural regime
for coordination-free convergence in distributed systems, alongside operation
commutativity (CRDTs) and invariant confluence (I-confluence).  The regime is
characterized by two conditions: well-founded compensation (WFC), which
ensures termination, and compensation commutativity (CC), which ensures
confluence.  Under both, the governance rewrite system has unique normal forms
by Newman's Lemma, and all processors consuming the same events converge to
the same valid state.  WFC and CC are sufficient for convergence; a uniform
bound on compensation depth (UBC) is orthogonal to convergence and governs
uniform complexity bounds.  Without UBC, compensation depth is unbounded;
without CC, different orders can reach different valid states.  CC as stated
over all states is sufficient but not necessary; its restriction to the states
reachable from the start is exactly necessary and sufficient when every
buffered event may fire and repair is canonical.

The result extends across organizational boundaries, in two layers.  Repair
confluence composes freely: on any acyclic network whose morphisms preserve
validity (M1), with a source-determined (R1), validity-preserving (R2)
\emph{resolution operator} at each multi-source target, one round of federated
normalization reaches the unique repair normal form from any state.  The core
mechanism is the authority argument: the sources' states determine every
target's shared component.  Event confluence does not compose for free.
Version~1 of this paper claimed that federated convergence requires only
validity preservation of the morphisms; that is false, because a target event
can read a shared value that morphism repair overwrites.  When events are
applied to normalized states, two local checks per edge suffice:
cross-registry compensation commutativity (C1) and repaired compensation
commutativity (C2).  The same two equations at the witnesses a run can reach
are exactly necessary and sufficient, and dropping either one admits
counterexamples that satisfy every component condition.  When the shared
domains are ordered lattices and repair is monotone, even \emph{cyclic}
networks have a unique repair normal form, the least fixed point (federally
valid when the network is valid at bottom), which reveals acyclicity as a
requirement of non-monotone repair only and exhibits order-independent repair as
chaotic iteration.  Event interleavings on a monotone cycle converge exactly when
independent events commute after re-normalization, which per-target C1 and C2
over the repair images imply.  Federated convergence is moreover
\emph{compositional} (proved on paper): a convex, internally convergent
sub-federation collapses to a single effective registry, so large networks
verify modularly and hierarchically.  Except for that collapse, the federated
results, their counterexamples and their exact converses are mechanized in
Coq/Rocq, axiom-free.

To make CC practically verifiable, we developed a verification calculus that
reduces the global CC obligation to local checks.  Invariant footprints
identify which invariants each event can affect; decomposable repair ensures
per-invariant fixes commute; and product composition enables modular
verification.  Federated convergence generalizes product composition from
independent registries to registries connected by directed semantic
constraints, with C1 and C2 as the per-edge obligations that the connection
adds.

Compensation commutativity is strictly weaker than the operation commutativity
required by CRDTs: operations need not commute, only their compensated results.
The gap between these two conditions is precisely the design space that
normalization confluence occupies---systems where operations are
non-commutative, invariant-violating, but convergent under repair.

The three regimes partition the coordination-free design space by where
convergence is guaranteed: in the operations, in the invariants, or in the
normalization rewrite system.  This last regime formalizes a pattern that
practitioners use implicitly---post-application repair as a convergence
mechanism---and provides the structural conditions under which it is sound, a
calculus for verifying those conditions, and a federated extension for
applying them across organizational boundaries.

\section*{Acknowledgments}

I thank my wife, Yuval, for her patience and support throughout this research.

\begin{thebibliography}{99}

\bibitem{blackwell2026nc}
Dayna~Blackwell,
``Normalization Confluence for Registry-Governed Stream Processing,''
Zenodo, 2026.
\textsc{doi}: \href{https://doi.org/10.5281/zenodo.18671870}{10.5281/zenodo.18671870}.

\bibitem{newman1942}
M.~H.~A. Newman,
``On theories with a combinatorial definition of `equivalence',''
\emph{Annals of Mathematics}, vol.~43, no.~2, pp.~223--243, 1942.

\bibitem{shapiro2011crdt}
M.~Shapiro, N.~Pregui\c{c}a, C.~Baquero, and M.~Zawirski,
``A comprehensive study of convergent and commutative replicated data types,''
\emph{INRIA Research Report RR-7506}, 2011.

\bibitem{hellerstein2010declarative}
J.~M.~Hellerstein,
``The declarative imperative: experiences and conjectures in distributed logic,''
\emph{SIGMOD Record}, vol.~39, no.~1, pp.~5--19, 2010.

\bibitem{ameloot2013relational}
T.~J.~Ameloot, F.~Neven, and J.~Van~den~Bussche,
``Relational transducers for declarative networking,''
\emph{Journal of the ACM}, vol.~60, no.~2, pp.~1--38, 2013.

\bibitem{bailis2014coordination}
P.~Bailis, A.~Fekete, M.~J.~Franklin, A.~Ghodsi, J.~M.~Hellerstein, and I.~Stoica,
``Coordination avoidance in database systems,''
\emph{Proc.\ VLDB Endowment}, vol.~8, no.~3, pp.~185--196, 2014.

\bibitem{baader1998term}
F.~Baader and T.~Nipkow,
\emph{Term Rewriting and All That},
Cambridge University Press, 1998.

\bibitem{vogels2009eventually}
W.~Vogels,
``Eventually consistent,''
\emph{Communications of the ACM}, vol.~52, no.~1, pp.~40--44, 2009.

\bibitem{garcia1987sagas}
H.~Garcia-Molina and K.~Salem,
``Sagas,''
\emph{Proc.\ ACM SIGMOD International Conference on Management of Data}, pp.~249--259, 1987.

\bibitem{lee1987synchronous}
E.~A.~Lee and D.~G.~Messerschmitt,
``Synchronous data flow,''
\emph{Proceedings of the IEEE}, vol.~75, no.~9, pp.~1235--1245, 1987.

\bibitem{kahn1974semantics}
G.~Kahn,
``The semantics of a simple language for parallel programming,''
\emph{Information Processing}, pp.~471--475, 1974.

\bibitem{akidau2015dataflow}
T.~Akidau, R.~Bradshaw, C.~Chambers, et~al.,
``The Dataflow Model: A practical approach to balancing correctness, latency, and cost in massive-scale, unbounded, out-of-order data processing,''
\emph{Proc.\ VLDB Endowment}, vol.~8, no.~12, pp.~1792--1803, 2015.

\bibitem{lamport1978time}
L.~Lamport,
``Time, clocks, and the ordering of events in a distributed system,''
\emph{Communications of the ACM}, vol.~21, no.~7, pp.~558--565, 1978.

\bibitem{tarski1955}
A.~Tarski,
``A lattice-theoretical fixpoint theorem and its applications,''
\emph{Pacific Journal of Mathematics}, vol.~5, no.~2, pp.~285--309, 1955.

\bibitem{cousot1977}
P.~Cousot and R.~Cousot,
``Abstract interpretation: a unified lattice model for static analysis of programs by construction or approximation of fixpoints,''
\emph{Proc.\ 4th ACM SIGACT-SIGPLAN Symp.\ on Principles of Programming Languages (POPL)}, pp.~238--252, 1977.

\bibitem{balegas2015indigo}
V.~Balegas, S.~Duarte, C.~Ferreira, R.~Rodrigues, N.~Pregui\c{c}a, M.~Najafzadeh, and M.~Shapiro,
``Putting consistency back into eventual consistency,''
\emph{Proc.\ 10th European Conf.\ on Computer Systems (EuroSys)}, 2015.

\bibitem{balegas2019ipa}
V.~Balegas, S.~Duarte, C.~Ferreira, R.~Rodrigues, and N.~Pregui\c{c}a,
``IPA: Invariant-preserving applications for weakly consistent replicated databases,''
\emph{Proc.\ VLDB Endowment}, vol.~12, no.~4, 2018.

\bibitem{meseguer1992}
J.~Meseguer,
``Conditional rewriting logic as a unified model of concurrency,''
\emph{Theoretical Computer Science}, vol.~96, no.~1, pp.~73--155, 1992.

\bibitem{duran2012coherence}
F.~Dur\'an and J.~Meseguer,
``On the Church-Rosser and coherence properties of conditional order-sorted rewrite theories,''
\emph{Journal of Logic and Algebraic Programming}, vol.~81, no.~7--8, pp.~816--850, 2012.

\bibitem{dijkstra1974}
E.~W.~Dijkstra,
``Self-stabilizing systems in spite of distributed control,''
\emph{Communications of the ACM}, vol.~17, no.~11, pp.~643--644, 1974.

\bibitem{gomes2017verifying}
V.~B.~F.~Gomes, M.~Kleppmann, D.~P.~Mulligan, and A.~R.~Beresford,
``Verifying strong eventual consistency in distributed systems,''
\emph{Proc.\ ACM on Programming Languages}, vol.~1, no.~OOPSLA, art.~109, 2017.

\end{thebibliography}

\appendix
\section{Errata and corrections}
\label{app:errata}

This appendix lists every statement of version~1 that version~2 refutes,
corrects or makes precise.  For each: the old statement (short), the
counterexample with its Coq name, and the corrected result with its Coq names.
All Coq names are theorems in the \texttt{Print Assumptions} gate of
\texttt{coq/verify.sh}, closed under the global context.  The body of the paper
states only the corrected results.

\subsection{Federated convergence}

\begin{description}[leftmargin=0pt,style=nextline]
\item[The federated rewrite system (\Cref{def:fed-comp,def:fed-grs}).]
  \emph{Old:} $\mathcal{G}_\Fed$ had apply steps ``as in the single-registry
  case'', enabled at any state.  \emph{Change:} version~2 adopts the guarded
  system (apply steps only at federally valid states), the semantics of a
  federated processor that applies each event to a normalized state, under
  which C1 and C2 are exact.  The unguarded system is treated in
  \Cref{rem:unguarded}: there C1 and C2 do not suffice
  (\texttt{fed\_grs\_c1\_c2\_insufficient}), and XU does
  (\texttt{fed\_thm\_fed\_convergence\_corrected}).

\item[\Cref{lem:authority} (Authority Determination).]
  \emph{Old (c):} the local component of a non-source ``is unconstrained by
  morphisms and converges via component CC''.  \emph{Counterexample:} a
  two-registry tree meeting every hypothesis in which two event orders leave
  the target's local counter at $0$ and at $1$ (\texttt{audit\_counterexample};
  against the exact hypotheses \texttt{fed\_lem\_authority\_c\_refuted}).
  \emph{Corrected:} (c) states only that repair never modifies local components
  (\texttt{fed\_lem\_authority\_c\_local}); the local normal form is determined
  under C1 and C2 by \Cref{thm:fed-convergence} (under XU in the unguarded
  system, \texttt{fed\_lem\_authority\_c\_corrected}).  (a) is stated precisely
  (a source's component of any normal form is the unique normal form of its own
  single-registry system on its own events): \texttt{fed\_lem\_authority\_a};
  (b): \texttt{fed\_lem\_authority\_b}.

\item[\Cref{lem:fed-termination,lem:resolved-termination} (Single-Round Termination).]
  \emph{Old:} one application of $\rho_\Fed$ yields a federally valid state.
  \emph{Change (strengthened):} from \emph{every} state, with the Phase~1 bound
  $\Phi_j(\sigma_j)$: \texttt{fed\_lem\_fed\_termination},
  \texttt{fed\_lem\_resolved\_termination}, \texttt{fed\_phase1\_bound}.  M1 is
  needed: \texttt{m1\_single\_round\_fails}.

\item[\Cref{thm:fed-cc} (Federated CC).]
  \emph{Old:} a tree-shaped network whose components satisfy WFC and CC
  satisfies federated CC.  \emph{Counterexamples:} \texttt{audit\_counterexample}
  (C1 fails) and \texttt{c2\_counterexample} (C1 holds, C2 fails); against the
  exact hypotheses \texttt{fed\_thm\_fed\_cc\_refuted}.  \emph{Corrected:} under
  (H), C1 and C2, the governed step satisfies CC1 at every federally valid state
  (\texttt{fed\_thm\_fed\_cc\_corrected\_machine}, \texttt{fed\_events\_commute});
  under XU, CC1 and CC2 hold at every state
  (\texttt{fed\_thm\_fed\_cc\_corrected}).

\item[\Cref{thm:fed-convergence} (Federated Normalization Confluence).]
  \emph{Old:} tree, M1, component WFC and CC imply that processors consuming
  the same events converge.  \emph{Counterexamples:}
  \texttt{fed\_thm\_fed\_convergence\_refuted} (C1 fails) and
  \texttt{fed\_c2\_paper\_counterexample} (C2 fails).  \emph{Corrected:} with C1
  and C2, the guarded system has unique, federally valid normal forms
  (\texttt{fed\_thm\_fed\_convergence\_guarded}); exactly, iff C1 and C2 hold at
  the reachable witnesses (\texttt{fed\_thm\_fed\_convergence\_exact},
  \texttt{fed\_guarded\_exact}, \texttt{fed\_exact}, \texttt{fed\_exact\_full},
  \texttt{acyclic\_gc\_iff}).  Non-vacuity:
  \texttt{fed\_supply\_paper\_instance}, \texttt{ms\_converges}.

\item[\Cref{cor:fed-nf} (Constructive Normal Form).]
  \emph{Old:} compute sources, propagate shared components, ``let each
  non-source converge its local component via component CC''.
  \emph{Counterexample:} even under C1 and C2 the recipe ``run the target's
  events as a single registry, then overwrite its shared component'' gives a
  state that no event order reaches (\texttt{fed\_cor\_fed\_nf\_refuted}).
  \emph{Corrected:} the normal form is the unique solution of the equations of
  \Cref{cor:fed-nf}, with repair interleaved after every target event
  (\texttt{fed\_cor\_fed\_nf\_corrected}, \texttt{fed\_nf\_constructive}); the
  old recipe holds under XU (\texttt{fed\_nf\_recipe\_xu}).

\item[\Cref{thm:resolved-convergence} and \Cref{cor:resolved-nf}.]
  \emph{Old:} the resolved analogues of \Cref{thm:fed-convergence} and
  \Cref{cor:fed-nf}, under M1, R1, R2 and component WFC and CC.
  \emph{Counterexample:} the tree counterexamples are resolved networks
  (\texttt{fed\_thm\_resolved\_convergence\_refuted}).  \emph{Corrected:}
  \texttt{fed\_thm\_resolved\_convergence\_guarded},
  \texttt{fed\_thm\_resolved\_convergence\_exact},
  \texttt{fed\_cor\_resolved\_nf\_corrected}; unguarded under XU,
  \texttt{fed\_thm\_resolved\_convergence\_corrected}.

\item[Remark ``Necessity of the resolution conditions''.]
  \emph{Old:} ``(R1) is necessary for well-definedness'' (of an unspecified
  object).  \emph{Made precise:} without (R1) the repair normal form is not
  unique: two federally valid normal forms with different shared components
  from one state (\texttt{r1\_necessary}).  The (R2) claim is mechanized
  (\texttt{r2\_necessary}).

\item[\Cref{thm:monotone-cycles} (Monotone Convergence Despite Cycles).]
  Three parts of the old statement fail.
  (i)~\emph{Old:} ``validity is preserved along the iteration by M1/R2, so
  $\bar{s}^\star$ is federally valid''.  \emph{Counterexample:} a 2-cycle of
  identity morphisms whose least fixed point is invalid
  (\texttt{fed\_thm\_monotone\_cycles\_lfp\_invalid}).  \emph{Corrected:} add
  validity at bottom (\texttt{net\_iter\_valid}, \texttt{net\_sweep\_valid},
  \texttt{net\_lfp\_valid}, \texttt{net\_Ncyc\_correct},
  \texttt{net\_Ncyc\_idem}).
  (ii)~\emph{Old:} $\bar{s}^\star = \bigsqcup_k \Phi^k(\bot)$ on any complete
  lattice.  \emph{Counterexample:} a monotone map on $0 < \cdots < \omega <
  \omega{+}1$ whose Kleene supremum is not fixed
  (\texttt{fed\_thm\_monotone\_cycles\_kleene\_formula\_fails}).
  \emph{Corrected:} the formula needs $\Phi$ to preserve the supremum of the
  Kleene chain (\texttt{kleene\_sup\_lfp}, \texttt{kleene\_sup\_valid});
  finite reachability needs ACC (\texttt{kleene\_acc\_lfp},
  \texttt{chaotic\_acc\_reaches\_lfp}).
  (iii)~\emph{Old:} ``consequently all processors consuming the same federated
  events converge''.  \emph{Counterexample:} a monotone 2-cycle meeting every
  hypothesis on which one allowed swap of two events diverges
  (\texttt{cyc\_counterexample},
  \texttt{fed\_thm\_monotone\_cycles\_events\_refuted}).  \emph{Corrected:}
  convergence iff GC (\texttt{cyc\_events\_converge\_iff},
  \texttt{net\_events\_converge\_iff}); per-target C1 and C2 over the repair
  images imply GC (\texttt{cyc\_check\_gc}, \texttt{cyc\_check\_gc\_lfp}).
  The source entry of $\Phi$ in \Cref{def:lattice-shared} now keeps the
  source's own shared value, as in the mechanized operator
  (\texttt{fed\_def\_lattice\_shared\_cyclic\_hyps}).

\item[\Cref{cor:acyclicity-monotone}.]
  \emph{Old:} acyclicity and monotonicity are two independent routes to
  ``federated confluence''.  \emph{Made precise:} to a unique repair normal
  form; event confluence needs C1 and C2 or GC in addition (counterexamples as
  for \Cref{thm:fed-convergence,thm:monotone-cycles};
  \texttt{negation\_not\_monotone}).

\item[\Cref{rem:infinite-lattice} (Infinite lattices and widening).]
  \emph{Old:} ``the convergence guarantee is unaffected'' without ACC, and
  widening yields a sound post-fixed point.  \emph{Made precise:} without ACC
  no iterate need reach the least fixed point even for continuous $\Phi$
  (\texttt{kleene\_sup\_instance}); a widened post-fixed point bounds the least
  fixed point (\texttt{widening\_sound}) but need not be a fixed point, so it
  is not the federated normal form (\texttt{widening\_not\_normal\_form}).

\item[\Cref{thm:collapse} (Compositional Collapse) and \Cref{cor:modular}.]
  \emph{Old:} (a) derived CC of the effective registry from internal
  convergence under component CC, via the refuted federated theorems.
  \emph{Corrected (on paper):} internal convergence now includes C1 and C2 on
  the internal edges (or the cyclic checks), and (a) gives CC1 of the governed
  step at valid states.  Not mechanized beyond the fold-append law
  (\texttt{rhoFold\_compositional}, \texttt{rhoF\_from\_app}).

\item[Remark ``Convexity and the monotone regime''.]
  \emph{Old:} if the sub-federation's repair is monotone then so is
  $\rho_\Fed^{\mathcal{J}}$.  \emph{Counterexample:}
  \texttt{fed\_rem\_convexity\_refuted} (Phase~1 is not monotone).
  \emph{Corrected:} monotone Phase~1 and $\Phi$ monotone in the locals as well
  give a monotone $\rho_\Fed^{\mathcal{J}}$ under ACC
  (\texttt{fed\_rem\_convexity\_corrected}, \texttt{lfp\_mono\_param}).

\item[Abstract, contributions 7 to 10, \Cref{rem:rewrite-property},
  \Cref{sec:regimes} and \Cref{tab:regimes}, the supply-chain application, and
  the conclusion.]
  \emph{Old:} ``federated convergence requires only validity preservation of
  the morphisms''; ``multi-registry systems inherit convergence from their
  components''; monotone repair gives convergence on any topology.
  \emph{Corrected:} repair confluence composes under M1/R1/R2 on acyclic
  networks (and under monotonicity with validity at bottom on cycles); event
  confluence needs C1 and C2 (acyclic) or GC (monotone cycles).  The
  counterexamples are those listed above.

\item[Examples.]
  \Cref{prop:cycle-necessary}, \Cref{prop:m1-necessary} and the
  manufacturer and supplier federation are unchanged and now mechanized
  (\texttt{prop\_cycle\_necessary}, \texttt{prop\_m1\_necessary},
  \texttt{ms\_instance}, \texttt{ms\_converges}).
\end{description}

%BASE-ERRATA

\end{document}
